\providecommand{\toplevelprefix}{../..}  %
\documentclass[../../book-main.tex]{subfiles}

\begin{document}

\chapter{Búsqueda de distribuciones de baja dimensión mediante reducción de ruido}
\label{ch:denoising}

\begin{quote}
    \hfill ``\textit{La información es la resolución de la incertidumbre}.''

    \hfill ---Claude Shannon, 1948
\end{quote}
\vspace{5mm}

En el \Cref{ch:classic-models}, mostramos cómo aprender clases simples de
distribuciones cuyos soportes se supone que son un único
subespacio de baja dimensión, una mezcla de tales subespacios o
gaussianas de bajo rango. Para simplificar, se considera que los diferentes modos lineales o gaussianos (ocultos)
son ortogonales o independientes,\footnote{O pueden
reducirse a tales casos idealizados en condiciones favorables.} tal y como
se ilustra en la \Cref{fig:subspaces}. Para estas distribuciones especiales,
se pueden derivar algoritmos de aprendizaje sencillos y eficaces con garantías de corrección y eficiencia,
siendo claras las interpretaciones geométricas y estadísticas
de las operaciones algorítmicas asociadas.

En la práctica, tanto la linealidad como la independencia son supuestos idealizados
que las distribuciones de los datos de alta dimensión del mundo real
\textit{rara vez} cumplen. Lo único que podemos suponer es que la
dimensión intrínseca de la distribución es muy baja en comparación con la
dimensión del espacio circundante en el que se encuentran los datos. Por lo tanto,
en este capítulo mostramos cómo aprender una clase más general de
distribuciones de baja dimensión en un espacio de alta dimensión que
\textit{no} es necesariamente (ni por tramos) lineal.

Las distribuciones del mundo real suelen contener múltiples componentes o
modos, por ejemplo, correspondientes a diferentes clases de objetos en
las imágenes. Es posible que estos modos no sean estadísticamente independientes 
e incluso puedan tener dimensiones intrínsecas diferentes. Por lo tanto, solemos
suponer que los datos se distribuyen en una mezcla de subvariedades no lineales 
de baja dimensión incrustadas en un espacio de alta dimensión. La 
\Cref{fig:mixture-manifolds} ilustra este tipo de distribución.  Otra
limitación es que, en situaciones reales, normalmente solo tenemos acceso
a un número finito de muestras de la distribución de datos,
en lugar de disponer de un conocimiento previo de la distribución exacta en sí.

Para aprender dicha distribución bajo estas condiciones, debemos abordar
tres preguntas fundamentales:
\begin{itemize}
    \item ¿Cuál es un enfoque general para aprender y representar una
        distribución de baja dimensión en un espacio de alta dimensión? 
    \item ¿Cómo medimos la complejidad de la distribución para
        poder aprovechar eficazmente su baja dimensionalidad en el aprendizaje? 
    \item ¿Cómo podemos hacer que el proceso de aprendizaje sea computacionalmente
        manejable y escalable, especialmente cuando la dimensión del entorno
        es elevada y el tamaño de la muestra es grande? 
\end{itemize}
Como veremos, la idea fundamental de la \textit{compresión}, o
\textit{reducción de dimensiones}, cuya eficacia ya hemos demostrado en el
caso lineal/independiente, sigue siendo un principio rector para el desarrollo
de modelos y métodos computacionales para distribuciones no lineales generales
de baja dimensión.

Dada su importancia teórica y práctica, analizamos con mayor
detalle cómo se materializa este marco basado en la compresión cuando la
distribución de interés puede modelarse o aproximarse adecuadamente mediante una
mezcla de subespacios de baja dimensión o gaussianas de bajo rango. Como se muestra en
los capítulos siguientes, estas clases de modelos, aparentemente restrictivas, sirven como
elementos básicos para algoritmos que realizan la compresión en distribuciones
generales.

\begin{figure}
    \centering
    \includegraphics[width=0.5\linewidth]{\toplevelprefix/chapters/chapter3/figs/mixed-manifolds.png}
    \caption{Los datos \(\vX\) se distribuyen en una mezcla de
        subvariedades de baja dimensión \(\cup_{j} \cM_{j}\) en un espacio
        envolvente de muy alta dimensión, por ejemplo, \(\R^{D}\). Este es un caso mucho
        más general que los casos especiales estudiados en el
        \Cref{ch:classic-models}, donde se suponía que las subvariedades eran
        lineales por tramos y, por lo tanto, admitían formas paramétricas y
    soluciones analíticas.}
    \label{fig:mixture-manifolds}
\end{figure}

\section{Minimización de la entropía y compresión}

\subsection{Entropía y tasas de codificación}
En el \Cref{ch:intro} mencionamos que el objetivo del aprendizaje es
encontrar la forma más sencilla de generar un conjunto de datos dado. En teoría,
la complejidad de Kolmogorov tenía como objetivo proporcionar una medida de
complejidad de este tipo, pero no es computable y no está asociada a ningún
esquema implementable que pueda reproducir realmente los datos. Por lo tanto,
necesitamos una medida de complejidad alternativa, computable y aplicable.
Esto nos lleva al concepto de \textit{entropía},
introducido por Shannon en 1948~\cite{Shannon-1948}.

Para ilustrar el carácter constructivo de la entropía, comencemos con
el caso más sencillo. Supongamos que tenemos una variable aleatoria discreta
que toma \(N\) valores distintos, o \textit{elementos}, \(\{\vx_{1},
\dots, \vx_{N}\}\) con una probabilidad igual de \(1/N\). Podemos
codificar cada token \(\vx_{i}\) utilizando la representación binaria de \(\log_{2} N\) bits
del número \(i\). Este esquema de codificación puede generalizarse para
codificar distribuciones discretas arbitrarias~\cite{Cover-Thomas}: dada
una distribución \(p\) tal que \(\sum_{i=1}^{N} p(\vx_{i}) = 1\), se
puede asignar a cada token \(\vx_{i}\) con probabilidad de \(p(\vx_{i})\)
un código binario de longitud \(\log_{2} [1/p(\vx_{i})] = - \log_{2}
p(\vx_{i})\) bits. Por lo tanto, el número promedio de bits, o \textit{la
tasa de codificación}, necesario para codificar cualquier muestra de la distribución
\(p(\cdot)\) viene dado por la expresión\footnote{Según la convención
    de la teoría de la información~\cite{Cover-Thomas}, el \(\log\) aquí tiene
base \(2\). Por lo tanto, la entropía se mide en bits (binarios).}:
\begin{equation}
    H(p) \doteq \Ex_{\vx \sim p}[-\log p(\vx)] = - \sum_{i=1}^{N}
    p(\vx_{i}) \log p(\vx_{i}).
    \label{eqn:entropy-discrete}
\end{equation}
Esto se conoce como la \textit{entropía} de la distribución (discreta)
\(p(\cdot)\). Tenga en cuenta que esta entropía es siempre no negativa y es
cero si y solo si \(p(\vx_{i}) = 1\) para algún \(\vx_{i}\) con \(i
\in [N]\).\footnote{Aquí utilizamos el hecho de que
\(\lim_{p \to 0} p \log p = 0\).} Para una variable aleatoria
\(\vx\), solemos escribir \(H(\vx)\) para referirnos a la entropía de la
distribución marginal de \(\vx\).

\subsection{Entropía diferencial}

Cuando la variable aleatoria \(\vx \in \R^{D}\) es continua y tiene una
densidad de probabilidad \(p\), se puede considerar que el límite de la suma anterior
\eqref{eqn:entropy-discrete} está relacionado con una integral:
\begin{equation}
    h(p) \doteq \Ex_{\vx \sim p}[-\log p(\vx)] = - \int_{\R^{D}}
    p(\vxi) \log p(\vxi) \odif{\vxi}.
    \label{eqn:entropy-differential}
\end{equation}
Más precisamente, dada una distribución de probabilidad continua \(p\) en
\(\R^{D}\), podemos discretizar \(\R^{D}\) en una unión de cubos disjuntos
\(C_{i}^{\epsilon}\) de longitud de arista \(\epsilon^{1/D}\), cada uno (digamos)
centrado en \(\vx_{i}\), y definir la distribución discretizada
\(p^{\epsilon}\) mediante \(p^{\epsilon}(\vx_{i}) = \Pr_{\vx \sim p}[\vx
\in C_{i}^{\epsilon}]\). Entonces se puede demostrar que \(H(p^{\epsilon}) +
\log \epsilon \approx h(p)\). Si tomamos \(\epsilon\) como un valor pequeño,
observamos que la entropía diferencial \(h(p)\) puede ser negativa.
Los lectores interesados pueden consultar \cite{Cover-Thomas} para obtener una explicación más
detallada. Al igual que con la entropía discreta, a menudo
escribiremos \(h(\vx)\) para referirnos a la entropía de la distribución (marginal) de \(\vx\).

\begin{example}[Entropía de las distribuciones gaussianas]
    Mediante un cálculo directo, se puede demostrar que la entropía de una
    variable aleatoria gaussiana \(x \sim \dNorm(\mu, \sigma^{2})\) es
    \begin{equation}
        h(x) = \frac{1}{2}\log(2\pi\sigma^{2}) + \frac{1}{2}.
        \label{eqn:entropy-Gaussian}
    \end{equation}
    La entropía de una variable aleatoria gaussiana multivariante \(\vx \sim
    \dNorm(\vmu, \vSigma)\) en \(\R^{D}\) es
    \begin{equation}
        h(\vx) = \frac{D}{2}(1 + \log(2\pi)) +
        \frac{1}{2}\logdet(\vSigma).
        \label{eqn:entropy-Gaussian-multi}
    \end{equation}
\end{example}

Al igual que ocurre con la entropía de una distribución discreta, nos gustaría que la
entropía diferencial estuviera relacionada con la tasa de codificación de algún
esquema de codificación viable. Por ejemplo, como se mencionó anteriormente, si
discretizamos el dominio de la distribución con una malla de tamaño
\(\epsilon > 0\) con centros \(\vx_{i}\), un posible esquema de codificación
consiste en asignar cada \(\vx\) a su centro más cercano \(\vx_{i}\) y
la entropía de la distribución discreta resultante se aproxima a la
entropía diferencial \cite{Cover-Thomas}.

Existen ciertas salvedades en la definición
de la entropía diferencial. En el caso de una distribución en un espacio de alta dimensión,
cuando su soporte se vuelve degenerado (de baja dimensión), su
entropía diferencial diverge hacia \(-\infty\). Este hecho se demuestra en el
\Cref{thm:max_entropy}, pero incluso en el caso explícito y sencillo de las
distribuciones gaussianas \eqref{eqn:entropy-Gaussian-multi}, cuando la
covarianza \(\vSigma\) es singular, podemos ver que
\(\logdet(\vSigma) = -\infty\), por lo que tenemos \(h(\vx) = -\infty\)
(Nota:
    Según el \Cref{thm:max_entropy}, las distribuciones gaussianas tienen la
mayor entropía entre todas las distribuciones con la misma media y covarianza). En
una situación así, no resulta evidente cómo cuantificar o
codificar adecuadamente dicha distribución. Sin embargo, las distribuciones (gaussianas)
degeneradas son precisamente los ejemplos más sencillos posibles, y podría decirse que los
más importantes, de distribuciones de baja dimensión en un espacio de alta dimensión.
En el próximo capítulo, analizaremos una solución completa
a esta aparente dificultad relacionada con la degeneración.

\subsection{Minimización de la tasa de codificación}
\label{sub:min_entropy}
El problema de aprendizaje consiste en recuperar una distribución (potencialmente continua)
\(p(\vx)\) a partir de muestras \(\{\vx_{1}, \dots, \vx_{N}\}\)
extraídas de ella. Para facilitar la exposición, escribimos \(\vX = \mat{\vx_{1},
\dots, \vx_{N}} \in \R^{D \times N}\). Dado que las distribuciones de
interés son (casi) de baja dimensión, esperamos que su entropía (diferencial)
sea muy pequeña. Sin embargo, a diferencia de las situaciones estudiadas en
el capítulo anterior, desconocemos la familia de modelos (analíticos)
de baja dimensión a la que pertenece \(p(\vx)\). Por lo tanto, comprobar
si la entropía es baja parece ser la única pauta en la que podemos
basarnos para identificar y modelar la distribución.

Ahora bien, si solo se tienen las muestras sin saber cuál es \(p(\vx)\), en
teoría podrían interpretarse como muestras de cualquier distribución
genérica. En particular, podrían interpretarse como cualquiera de las siguientes:
\begin{enumerate}
    \item como muestras de la propia distribución empírica \(p^{\vX}\),
        que asigna una probabilidad de \(1/N\) a cada una de las
        \(N\) muestras \(\vx_{i}, i \in [N]\); 
    \item como muestras de una distribución normal estándar \(\vx^{n}
        \sim p^{n} \doteq \dNorm(\vzero, \sigma^{2} \vI)\) con
        una varianza \(\sigma^{2}\) lo suficientemente grande (digamos, mayor que las
        normas al cuadrado de las muestras); 
    \item como muestras de una distribución normal \(\vx^{e} \sim p^{e}
        \doteq \dNorm(\vzero, \hat{\vSigma})\), siendo la covarianza
        \(\hat{\vSigma} = \frac{1}{N} \vX \vX^{\top}\) la
        covarianza empírica de las muestras; 
    \item como muestras de una distribución \(\hat{\vx} \sim
        \hat{q}(\vx)\) que se aproxima mucho a la distribución
        verdadera \(p\). 
\end{enumerate}
La pregunta es: ¿qué interpretación es mejor y en qué sentido? Supongamos
que crees que estos datos \(\vX\) provienen de una distribución concreta
\(q(\vx)\), que podría ser una de las anteriores. Entonces podríamos
codificar los puntos de datos con el libro de códigos óptimo para la distribución
\(q(\vx)\). Entonces, a medida que el número de muestras \(N\) aumenta, la
longitud media de codificación requerida (o tasa de codificación) viene dada por la
denominada \textit{entropía cruzada}:
\begin{equation}
    \CE(p \mmid q) \doteq \Ex_{\vx \sim p}[-\log q(\vx)].
\end{equation}
Si hemos identificado
la distribución correcta \(p(\vx)\), la tasa de codificación viene dada por la
entropía \(h(p) = \Ex_{\vx \sim p}[-\log p(\vx)]\). Resulta
que la longitud de codificación anterior \(\CE(p \mmid q)\) es siempre mayor
o igual que la entropía \(h(\vx)\), a menos que \(p = q\).

Por lo tanto, dado un conjunto de datos muestreados \(\vX\), para determinar qué
distribución caracteriza mejor los datos entre
\(p^{n}\), \(p^{e}\) y \(\hat{q}\), podemos comparar sus tasas de codificación
para \(\vX\) y ver cuál ofrece la tasa más baja. Sabemos
por lo anterior que la tasa de codificación (teóricamente alcanzable) de una
distribución está estrechamente relacionada con su entropía. En general,
\begin{equation}
    h(p^{n}) > h(p^{e}) > h(\hat{q}).
\end{equation}
Si los datos \(\vX\) se codificaran mediante el libro de códigos asociado
a cada una de estas distribuciones, la tasa de codificación de \(\vX\)
disminuiría, por lo tanto, en el mismo orden:
\begin{equation}
    \CE(p \mmid p^{n}) > \CE(p \mmid p^{e}) > \CE(p \mmid \hat{q}).
\end{equation}

Esta observación ofrece una pauta general para buscar una
distribución \(p(\vx)\) que tenga una estructura de baja dimensión.
Sugiere dos posibles enfoques:
\begin{enumerate}[label={A\arabic*.}, leftmargin=*]
    \item Partiendo de una distribución general (por ejemplo, una distribución
        normal) con alta entropía, transformar gradualmente la
        distribución hacia la distribución (empírica) de los datos
        mediante la reducción de la entropía; 
    \item Dentro de una amplia familia de distribuciones (paramétricas o no paramétricas)
        con esquemas de codificación explícitos que codifican los
        datos dados, buscar progresivamente mejores esquemas de codificación
        que ofrezcan tasas de codificación más bajas. 
\end{enumerate}
En teoría, tanto el enfoque A1 como el A2 pretenden lograr lo mismo.
Para A1, debemos asegurarnos de que exista tal ruta de transformación y de que sea
computable. Para A2, es necesario que la familia elegida sea lo suficientemente amplia
y pueda aproximarse bastante (o abarcar) la distribución real. En cualquiera de 
los dos enfoques, debemos asegurarnos de que las soluciones con menor entropía 
o mejores tasas de codificación puedan calcularse de manera eficiente
y converjan rápidamente hacia la distribución deseada.\footnote{Por ejemplo, 
la distribución de datos del mundo real, como imágenes y textos.} Exploraremos
ambos enfoques en las secciones restantes de este capítulo.

\section{Compresión mediante reducción de ruido}\label{sub:compression_denoising}

En esta sección describimos una forma \textit{natural} y
\textit{computacionalmente manejable} de aprender una distribución
\(p_{\vx}\) mediante el aprendizaje de una codificación paramétrica que minimice la
entropía o la tasa de codificación de la representación y luego utilizando esta codificación
para transformar muestras de alta entropía de una distribución gaussiana estándar en
muestras de baja entropía de la distribución objetivo, tal como se ilustra en
la \Cref{fig:diffusion-chapter3}. Esta metodología combina los enfoques
A1 y A2 mencionados anteriormente para aprender y obtener muestras de la distribución.

\begin{figure}[t]
    \centering
    \includegraphics[width=\linewidth]{\toplevelprefix/chapters/chapter3/figs/diffusion_pipeline.png}
    \caption{Ilustración de un proceso iterativo de eliminación de ruido que,
        partiendo de una distribución gaussiana isotrópica, converge hacia una
    distribución de datos arbitraria.}
    \label{fig:diffusion-chapter3}
\end{figure}

\subsection{Procesos de difusión y de reducción de ruido}\label{sub:intro_diffusion_denoising}

En primer lugar, queremos un procedimiento que reduzca la entropía de una variable 
aleatoria con mucho ruido, generando una muestra de menor entropía a partir de la 
distribución de los datos. Aquí describimos un posible enfoque, quizá el más natural.
Primero, diseñamos un método para \textit{aumentar gradualmente} la entropía de
las muestras ya extraídas de la distribución de datos. A continuación, buscamos una
\textit{inversa aproximada} de este proceso. En general, el aumento de
la entropía no tiene inverso, ya que parte de la información de la distribución original
puede destruirse. Por lo tanto, nos limitamos a un caso especial en el que (1) la 
operación que aumenta la entropía es simple, computable y reversible, así como (2) 
contamos con una codificación (paramétrica) de la distribución de los datos, tal como 
se ha mencionado anteriormente. Como veremos, estas dos
condiciones hacen que nuestro enfoque sea viable.

\subsubsection{Procesos de difusión}

Aumentaremos la entropía de la forma más sencilla posible,
es decir, \textit{añadiendo ruido gaussiano isotrópico}. Más concretamente, dada
la variable aleatoria \(\vx\), podemos considerar el \textit{proceso estocástico}
\((\vx_{t})_{t \in [0, T]}\) que le añade ruido gradualmente, es decir,
\begin{equation}\label{eq:additive_gaussian_noise_model}
    \vx_{t} \doteq \vx + t\vg, \qquad \forall t \in [0, T],
\end{equation}
donde \(T \in [0, \infty)\) es un horizonte temporal y \(\vg \sim
\dNorm(\vzero, \vI)\) se extrae independientemente de \(\vx\). Este proceso
es un ejemplo de \textit{proceso de difusión}, llamado así porque
distribuye (es decir, \textit{difunde}) la masa de probabilidad por todo
\(\R^{D}\) a medida que pasa el tiempo, aumentando la entropía. Esta
intuición queda confirmada gráficamente en la \Cref{fig:ve_forward_density},
así como de forma rigurosa mediante el siguiente teorema.
\begin{theorem}[Versión simplificada del \Cref{thm:diffusion_entropy_increases}]
    Supongamos que \((\vx_{t})_{t \in [0, T]}\) sigue el modelo
    \eqref{eq:additive_gaussian_noise_model}. Para cualquier \(t \in (0,
    T]\), la variable aleatoria \(\vx_{t}\) tiene una densidad de probabilidad
    \(p_{t}\) y una entropía diferencial
    \(h(\vx_{t}) > -\infty\). Además, bajo ciertas condiciones técnicas
    sobre \(\vx\),
    \begin{equation}
        \odv*{h(\vx_{t})}{t} > 0, \qquad \forall t \in (0, T],
    \end{equation}
    es decir, la entropía de la distribución de \(\vx_{t}\) aumenta
    con el paso del tiempo \(t\).
\end{theorem}
La demostración es elemental, pero bastante extensa, por lo que la posponemos hasta
la \Cref{sub:diffusion_entropy_increases}. La principal consecuencia, aún no expresada,
de este resultado es que \(h(\vx_{t}) > h(\vx_{s})\) para todo
\(t > s \geq 0\). Para comprobarlo, nótese en que si \(s = 0\) y \(h(\vx) =
h(\vx_{0}) = -\infty\), entonces \(h(\vx_{t}) > -\infty\) para todo \(t >
0\), por lo que la conclusión es verdadera; de lo contrario, \(h(\vx_{t}) = h(\vx_{s}) +
\int_{s}^{t} \bs{\odv*{h(\vx_{u})}{u}} \odif{u} > h(\vx_{s})\) por el
teorema fundamental del cálculo, por lo que en ambos casos \(h(\vx_{t}) >
h(\vx_{s})\) para todo \(t > s \geq 0\).

\begin{figure}
    \includegraphics[width=\textwidth]{\toplevelprefix/chapters/chapter3/figs/ve_forward_diffusion_density.png}
    \caption{\small\textbf{Difusión de una mezcla de distribuciones gaussianas \(\vx\).}
        De izquierda a derecha, observamos la evolución de la densidad
        \(p_{t}\) de \(\vx_{t}\) a medida que \(t\) varía de \(0\) a \(10\),
        junto con algunas muestras representativas (en rojo). El plano está
        coloreado según la densidad de probabilidad \(p_{t}\); las regiones de alta densidad
        aparecen en un azul más oscuro. Observamos que la masa de
        probabilidad se vuelve menos concentrada a medida que aumenta \(t\), lo que indica que
    la entropía aumenta.}
    \label{fig:ve_forward_density}
\end{figure}

\begin{example}[Añadir ruido a una mezcla de
    distribuciones gaussianas]\label{ex:gaussian_mixture}
    En este ejemplo, consideramos añadir ruido a una mezcla de \(K\)
    distribuciones gaussianas. Elegir \(K\) significa
    \(\vmu_{1}, \dots, \vmu_{K} \in \R^{D}\), \(K\) matrices de covarianza
    \(\vSigma_{1}, \dots, \vSigma_{K} \in \PSD(D)\) y
    probabilidades \(\pi_{1}, \dots, \pi_{K} \in [0, 1]\) tales que
    \(\sum_{k = 1}^{K}\pi_{k} = 1\). Entonces, \(\vx\) se muestrea de la
    siguiente manera.
    \begin{itemize}
        \item En primer lugar, se toma una muestra de un índice (o \textit{etiqueta}) \(y \in [K]\)
            de tal manera que \(y = k\) con probabilidad \(\pi_{k}\). 
        \item En segundo lugar, \(\vx\) se obtiene mediante muestreo de la distribución gaussiana
            \(\dNorm(\vmu_{y}, \vSigma_{y})\); es decir, si \(y = k\), entonces
            \(\vx \sim \dNorm(\vmu_{k}, \vSigma_{k})\). 
    \end{itemize}
    Por lo tanto, \(\vx\) sigue una distribución de mezcla gaussiana
    \begin{equation}
        \vx \sim \sum_{k = 1}^{K}\pi_{k}\dNorm(\vmu_{k}, \vSigma_{k}).
    \end{equation}
    Los modelos de mezcla gaussiana son relevantes porque pueden aproximar
    muchas distribuciones: dada cualquier distribución para \(\vx\), existe
    una mezcla gaussiana que la aproxima con una precisión arbitraria.
    Por este motivo, los modelos de mezcla gaussiana son una alternativa popular y
    analíticamente manejable para las distribuciones de datos multimodales y complejas.
    En este libro los utilizamos ampliamente como «problema modelo» para comprender 
    y desarrollar casi todos los conceptos. Por ahora, representamos los datos \(\vx\) 
    mediante una mezcla gaussiana. Tras añadir
    ruido, la distribución de \(\vx_{t}\) vuelve a ser una mezcla gaussiana:
    \begin{equation}
        \vx_{t} = \vx + t\vg \sim \sum_{k =
        1}^{K}\pi_{k}\dNorm(\vmu_{k}, \vSigma_{k} + t^{2}\vI).
    \end{equation}
    Para ver con un ejemplo cómo evoluciona la entropía, supongamos que \(K = 1\), de modo que
    \(\vx \sim \dNorm(\vmu, \vSigma)\). Según
    \eqref{eqn:entropy-Gaussian-multi}, la entropía de \(\vx_{t}\) es
    \begin{equation}
        h(\vx_{t}) = \frac{D}{2}(1 + \log(2\pi)) +
        \frac{1}{2}\logdet(\vSigma + t^{2}\vI).
    \end{equation}
    Al derivar con respecto al tiempo (véase el \Cref{app:optimization})
    se obtiene
    \begin{equation}
        \odv*{h(\vx_{t})}{t} = t\tr((\vSigma + t^{2}\vI)^{-1}) =
        \sum_{i = 1}^{D}\frac{t}{\lambda_{i}(\vSigma) + t^{2}} > 0,
    \end{equation}
    donde \(\lambda_{i}(\vSigma)\) son los valores propios (no negativos) de
    \(\vSigma\). Dado que esta derivada es positiva, la entropía
    \(h(\vx_{t})\) de la distribución con ruido aumenta con el tiempo.
\end{example}

\subsubsection{Procesos de reducción de ruido}
La operación inversa a la adición de ruido se conoce como \textit{eliminación o reducción de ruido}.
Se trata de un tema clásico y ampliamente estudiado en el procesamiento de señales y
la teoría de sistemas, que se materializa (por ejemplo) en el filtro de Wiener y
el filtro de Kalman (véase la \Cref{sec:analytical}). Varios problemas tratados en
el \Cref{ch:classic-models}, como el PCA, el ICA y el aprendizaje de diccionarios, son
casos concretos del problema de la eliminación de ruido. Para un \(t\) fijo y
el modelo de ruido gaussiano aditivo
\eqref{eq:additive_gaussian_noise_model}, el problema de eliminación de ruido puede
formularse como el intento de aprender una función \(\bar{\vx}^{\star}(t,
\cdot)\) que constituya la mejor aproximación posible (en esperanza)
de la variable aleatoria verdadera \(\vx\), dados tanto \(t\) como \(\vx_{t}\):
\begin{equation}\label{eq:denoising_loss}
    \bar{\vx}^{\star}(t, \cdot) \in \argmin_{\bar{\vx}(t,
    \cdot)} \Ex \norm{\vx - \bar{\vx}(t, \vx_{t})}_{2}^{2}.
\end{equation}
La solución a este problema, al optimizar \(\bar{\vx}(t, \cdot)\)
sobre todas las funciones posibles (integrables en el cuadrado), es el así llamado
\textit{filtro óptimo de reducción de ruido de Bayes} y resulta ser la esperanza condicional:
\begin{equation}\label{eq:optimal_denoiser}
    \bar{\vx}^{\star}(t, \vxi) \doteq \Ex[\vx \mid \vx_{t} = \vxi].
\end{equation}

Esta expresión justifica la notación \(\bar{\vx}\), que calcula
una esperanza condicional (es decir, la media condicional o el promedio
condicional). En resumen, intenta eliminar el ruido de la entrada ruidosa
y genera la mejor estimación posible (en términos de esperanza y con
respecto a la distancia \(\ell^{2}\)) de la variable aleatoria original sin ruido.

\begin{example}[Eliminación del ruido gaussiano a partir de una mezcla de
    distribuciones gaussianas]\label{example:denoising_gaussian_mixture}
    En este ejemplo calculamos el filtro óptimo de reducción de ruido de Bayes para una
    clase importante de distribuciones: el modelo de mezcla gaussiana
    del \Cref{ex:gaussian_mixture}. Definamos \(\phi(\vx ;
    \vmu, \vSigma)\) como la densidad de probabilidad de \(\dNorm(\vmu,
    \vSigma)\) evaluada en \(\vx\). Con esta notación, la densidad de
    \(\vx_{t}\) es
    \begin{equation}
        p_{t}(\vx_{t}) = \sum_{k = 1}^{K}\pi_{k}\phi(\vx_{t} ;
        \vmu_{k}, \vSigma_{k} + t^{2}\vI).
    \end{equation}

    Condicionadas en \(y\), las variables siguen una distribución gaussiana conjunta: si
    escribimos \(\vx = \vmu_{y} + \vSigma_{y}^{1/2}\vu\), donde
    \((\cdot)^{1/2}\) es la raíz cuadrada matricial y \(\vu \sim
    \dNorm(\vzero, \vI)\) independientemente de \(y\) (y de \(\vg\)), entonces tenemos
    \begin{equation}
        \mat{\vx \\ \vx_{t}} = \mat{\vmu_{y} \\ \vmu_{y}} +
        \mat{\vSigma_{y}^{1/2} & \vzero \\ \vSigma_{y}^{1/2} &
        t\vI}\mat{\vu \\ \vg}.
    \end{equation}
    Esto muestra que \(\vx\) y \(\vx_{t}\) siguen una distribución gaussiana conjunta
    (condicionada en \(y\)), tal como se afirmaba. Por lo tanto, podemos escribir
    \begin{equation}
        \mat{\vx \\ \vx_{t}} \sim \dNorm\rp{\mat{\vmu_{y}
            \\ \vmu_{y}}, \mat{\vSigma_{y} & \vSigma_{y} \\ \vSigma_{y} &
        \vSigma_{y} + t^{2}\vI}}.
    \end{equation}
    Por lo tanto, la esperanza condicional de \(\vx\) dada \(\vx_{t}\) (es decir, el
    filtro óptimo de reducción de ruido de Bayes condicionado en \(y\)) es, como es bien sabido,
    (véase también el \Cref{exercise:conditional_gaussian})
    \begin{equation}
        \Ex[\vx \mid \vx_{t}, y] = \vmu_{y} + \vSigma_{y}(\vSigma_{y}
        + t^{2}\vI)^{-1}(\vx_{t} - \vmu_{y}).
    \end{equation}
    Para hallar el filtro de ruido óptimo de Bayes global, utilizamos la ley de
    la esperanza iterada, obteniendo
    \begin{align}
        \bar{\vx}^{\star}(t, \vx_{t})
        &= \Ex[\vx \mid \vx_{t}] \\
        &= \Ex[\Ex[\vx \mid \vx_{t}, y] \mid \vx_{t}] \\
        &= \sum_{k = 1}^{K}\Pr[y = k \mid \vx_{t}]\Ex[\vx \mid \vx_{t}, y = k].
    \end{align}
    La probabilidad se puede tratar de la siguiente manera. Sea \(p_{t \mid
    y}\) la densidad de probabilidad de \(\vx_{t}\) condicionada al
    valor de \(y\). Entonces
    \begin{align}
        \Pr[y = k \mid \vx_{t}]
        &= \frac{p_{t \mid y}(\vx_{t} \mid k)\pi_{k}}{p_{t}(\vx_{t})} \\
        &= \frac{\pi_{k}\phi(\vx_{t} ; \vmu_{k}, \vSigma_{k}
        + t^{2}\vI)}{\sum_{i = 1}^{K}\pi_{i}\phi(\vx_{t} ; \vmu_{i},
        \vSigma_{i} + t^{2}\vI)}.
    \end{align}
    Por otro lado, la esperanza condicional es la descrita anteriormente:
    \begin{equation}
        \Ex[\vx \mid \vx_{t}, y = k] = \vmu_{k} +
        \vSigma_{k}(\vSigma_{k} + t^{2}\vI)^{-1}(\vx_{t} - \vmu_{k}).
    \end{equation}
    En resumen, el verdadero filtro de eliminación de ruido óptimo de Bayes es
    \begin{align}\label{eq:gmm_bayes_optimal_denoiser}
        \bar{\vx}^{\star}(t, \vx_{t}) &= \sum_{k =
        1}^{K}\Bigg\{\frac{\pi_{k}\phi(\vx_{t}
            ;\vmu_{k}, \vSigma_{k} + t^{2}\vI)}{\sum_{i =
                1}^{K}\pi_{i}\phi(\vx_{t}
            ;\vmu_{i}, \vSigma_{i} + t^{2}\vI)}\\
            &\qquad \qquad \qquad \qquad   \cdot \bp{\vmu_{k} +
        \vSigma_{k}(\vSigma_{k} + t^{2}\vI)^{-1}(\vx_{t} - \vmu_{k})}\Bigg\}.
    \end{align}
    Este ejemplo es especialmente importante, y varios casos particulares
    nos proporcionarán una gran comprensión conceptual más adelante. Por ahora, intentemos
    extraer alguna intuición geométrica a partir de la forma funcional
    del filtro óptimo de reducción de ruido \eqref{eq:gmm_bayes_optimal_denoiser}.

    Para comprender \eqref{eq:gmm_bayes_optimal_denoiser}
    de manera intuitiva, fijemos primero \(K = 1\) (es decir, una sola gaussiana) de tal
    manera que \(\vx \sim \dNorm(\vmu, \vSigma)\). Diagonalicemos entonces
    \(\vSigma = \vV\vLambda \vV^{\top}\). Entonces, el filtro óptimo de eliminación de ruido de Bayes es
    \begin{align}
        \bar{\vx}^{\star}(t, \vx_{t}) & = \vmu + \vSigma(\vSigma +
        t^{2}\vI)^{-1}(\vx_{t} - \vmu) \\
        & = \vmu + \vV\mat{\lambda_{1}/(\lambda_{1} + t^{2}) & & \\ &
            \ddots & \\ & & \lambda_{D}/(\lambda_{D} +
        t^{2})}\vV^{\top}(\vx_{t} - \vmu),\label{eq:gaussian_denoiser}
    \end{align}
    donde \(\lambda_{1}, \dots, \lambda_{D}\) son los valores propios de
    \(\vSigma\). Podemos observar que este filtro de eliminación de ruido consta de tres pasos:
    \begin{itemize}
        \item Trasladar la entrada \(\vx_{t}\) en \(\vmu\). 
        \item Contraer la entrada (trasladada) \(\vx_{t} - \vmu\) en
            cada dirección de vector propio en una cantidad
            \(\lambda_{i}/(\lambda_{i} + t^{2})\). Si la entrada trasladada
            es de rango bajo y algunos valores propios de \(\vSigma\) son
            cero, el filtro de reducción de ruido contrae inmediatamente estas direcciones hasta
            \(0\), lo que garantiza que la salida de la
            contracción sea igualmente de bajo rango.
        \item Desplazar el resultado hacia atrás en \(\vmu\). 
    \end{itemize}
    Esta es la interpretación geométrica del filtro de reducción de ruido de una
    \textit{única} gaussiana. El filtro de ruido global del modelo de mezcla gaussiana
    \eqref{eq:gmm_bayes_optimal_denoiser} utiliza \(K\)
    filtros de ruido de este tipo, ponderando su salida mediante las probabilidades a posteriori
    \(\Pr[y = k \mid \vx_{t}]\). Si las medias de las
    distribuciones gaussianas están bien separadas, estas probabilidades a posteriori son
    muy cercanas a \(0\) o \(1\) cerca de cada media o grupo. En este
    régimen, el filtro de ruido global
    \eqref{eq:gmm_bayes_optimal_denoiser} tiene la misma interpretación geométrica
    que el filtro de ruido gaussiano único mencionado anteriormente.
\end{example}

\begin{figure}
    \centering
    \includegraphics[width=\textwidth]{\toplevelprefix/chapters/chapter3/figs/ve_forward_diffusion_denoising.png}\vspace{-0.15in}
    \caption{\small \textbf{Filtro de eliminación de ruido óptimo de Bayes y puntuación de un
        modelo de mezcla gaussiana.} En el mismo contexto que en la
        \Cref{fig:ve_forward_density}, mostramos el efecto del
        filtro de ruido óptimo de Bayes \(\bar{\vx}^{\star}\) representando gráficamente
        \(\vx_{t}\) (rojo) y \(\bar{\vx}^{\star}(t, \vx_{t})\) (verde)
        para una selección de \(t\) y \(\vx_{t}\). Según la fórmula de Tweedie
        \Cref{thm:tweedie}, el residuo entre ambos es proporcional a
        la denominada puntuación (Hyv\"arinen) \(\nabla_{\vx_{t}}\log
        p_{t}(\vx_{t})\). Podemos observar que la puntuación apunta hacia los
    modos de la distribución de \(\vx_{t}\).}
    \label{fig:ve_forward_denoising}
\end{figure}

Intuitivamente y como podemos ver en el
\Cref{example:denoising_gaussian_mixture}, el filtro de ruido óptimo de Bayes
\(\bar{\vx}^{\star}(t, \cdot)\) debería desplazar su entrada \(\vx_{t}\)
hacia los modos de la distribución de \(\vx\). Resulta que podemos
cuantificar esto mostrando que el filtro de ruido óptimo de Bayes
\textit{da un paso de ascenso por gradiente} en la densidad (logarítmica)
de \(\vx_{t}\), que (recordemos) hemos denotado como \(p_{t}\). Es decir,
aplicar el filtro de eliminación de ruido implica pasar de la iteración de entrada a una
región de mayor probabilidad dentro de esta distribución (perturbada). Para valores 
pequeños de \(t\), la perturbación es pequeña, por lo que nuestra intuición inicial 
es, por lo tanto, casi totalmente correcta. La imagen se muestra en la
\Cref{fig:ve_forward_denoising} y se formula rigurosamente como
la fórmula de Tweedie \cite{Robbins1956AnEB}.
\begin{theorem}[La fórmula de Tweedie]\label{thm:tweedie}
    Supongamos que \((\vx_{t})_{t \in [0, T]}\) obedece a
    \eqref{eq:additive_gaussian_noise_model}. Sea \(p_{t}\) la
    densidad de \(\vx_{t}\) (tal como se ha definido anteriormente). Entonces
    \begin{equation}\label{eq:tweedie}
        \Ex[\vx \mid \vx_{t}] = \vx_{t} + t^{2}\nabla_{\vx_{t}} \log
        p_{t}(\vx_{t}).
    \end{equation}
\end{theorem}
\begin{proof}
    Para demostrarlo, supongamos que \(\vx\) tiene una
    densidad,\footnote{Obsérvese que el
        teorema es cierto sin esta suposición y que puede demostrarse mediante
        una modificación mínima de esta demostración para utilizar medidas de probabilidad
    en lugar de densidades.} y llamemos a esta densidad \(p\). Sean
    \(p_{0 \mid t}\) y \(p_{t \mid 0}\) las densidades condicionales de \(\vx
    = \vx_{0}\) dado \(\vx_{t}\) y de \(\vx_{t}\) dado \(\vx\), respectivamente.
    Sea \(\phi(\vx ; \vmu, \vSigma)\) la densidad de \(\dNorm(\vmu,
    \vSigma)\) evaluada en \(\vx\), de modo que \(p_{t \mid 0}(\vx_{t} \mid \vx)
    = \phi(\vx_{t} ; \vx, t^{2}\vI)\). Entonces, un cálculo sencillo da
    \begin{align}
        \nabla_{\vx_{t}} \log p_{t}(\vx_{t})
        &= \frac{\nabla_{\vx_{t}}p_{t}(\vx_{t})}{p_{t}(\vx_{t})} \\
        &=
        \frac{1}{p_{t}(\vx_{t})}\nabla_{\vx_{t}}\int_{\R^{D}}p(\vx)p_{t
        \mid 0}(\vx_{t} \mid \vx)\odif{\vx} \\
        &=
        \frac{1}{p_{t}(\vx_{t})}\nabla_{\vx_{t}}\int_{\R^{D}}p(\vx)\phi(\vx_{t}
        ;\vx, t^{2}\vI)\odif{\vx} \\
        &=
        \frac{1}{p_{t}(\vx_{t})}\int_{\R^{D}}p(\vx)[\nabla_{\vx_{t}}\phi(\vx_{t}
        ;\vx, t^{2}\vI)]\odif{\vx} \\
        &= \frac{1}{p_{t}(\vx_{t})}\int_{\R^{D}}p(\vx)\phi(\vx_{t} ;
        \vx, t^{2}\vI)\bs{-\frac{\vx_{t} - \vx}{t^{2}}}\odif{\vx} \\
        &=
        \frac{1}{t^{2}p_{t}(\vx_{t})}\int_{\R^{D}}p(\vx)\phi(\vx_{t}
        ; \vx, t^{2}\vI)[\vx - \vx_{t}]\odif{\vx} \\
        &= \frac{1}{t^{2}p_{t}(\vx_{t})}\int_{\R^{D}}p(\vx)\phi(\vx_{t} ; \vx,
        t^{2}\vI)\vx \odif{\vx} \\
        & \quad -
        \frac{\vx_{t}}{t^{2}p_{t}(\vx_{t})}\int_{\R^{D}}p(\vx)\phi(\vx_{t}
        ; \vx, t^{2}\vI)\odif{\vx}\nonumber \\
        &= \frac{1}{t^{2}p_{t}(\vx_{t})}\int_{\R^{D}}p(\vx)p_{t \mid
        0}(\vx_{t} \mid \vx)\vx \odif{\vx} -
        \frac{\vx_{t}}{t^{2}p_{t}(\vx_{t})}p_{t}(\vx_{t}) \\
        &=
        \frac{1}{t^{2}p_{t}(\vx_{t})}\int_{\R^{D}}p_{t}(\vx_{t})p_{0
        \mid t}(\vx \mid \vx_{t})\vx \odif{\vx} -
        \frac{\vx_{t}}{t^{2}p_{t}(\vx_{t})}p_{t}(\vx_{t}) \\
        &= \frac{1}{t^{2}}\int_{\R^{D}}p_{0 \mid t}(\vx \mid
        \vx_{t})\vx \odif{\vx} - \frac{\vx_{t}}{t^{2}} \\
        &= \frac{1}{t^{2}}\Ex[\vx \mid \vx_{t}] - \frac{\vx_{t}}{t^{2}} \\
        &= \frac{\Ex[\vx \mid \vx_{t}] - \vx_{t}}{t^{2}}.
    \end{align}
    Una simple reorganización de la igualdad anterior demuestra el teorema.
\end{proof}
Este resultado establece una conexión entre la eliminación de ruido y la optimización:
el filtro de ruido óptimo de Bayes da un único paso de ascenso por gradiente en
la densidad de datos perturbada \(p_{t}\), y el tamaño del paso se reduce de forma adaptativa
(es decir, da pasos más precisos) a medida que la perturbación
de la distribución de datos se va reduciendo. La cantidad
\(\nabla_{\vx_{t}}\log p_{t}(\vx_{t})\) se denomina
\textit{puntuación (de Hyv\"arinen)} y aparece con frecuencia en debates
sobre la eliminación de ruido, etc.; apareció por primera vez en un artículo de Aapo
Hyv\"arinen en el contexto del ICA \cite{hyvarinen05a}.

De manera similar a cómo un solo paso de descenso por gradiente casi nunca es suficiente para
minimizar un objetivo en la práctica cuando se inicia lejos del óptimo, la
salida del filtro de ruido óptimo de Bayes \(\bar{\vx}^{\star}(t, \cdot)\)
casi nunca se encuentra en una región de alta probabilidad de la distribución de datos
cuando \(t\) es grande, \textit{especialmente} cuando los datos
tienen una estructura de baja dimensión. Ilustramos este punto de manera explícita
en el siguiente ejemplo.
\begin{example}[Eliminación de ruido en una mezcla de
    dos puntos]\label{example:denoising_twopoints}
    Sea \(x\) una variable aleatoria uniforme en el conjunto de dos puntos \(\{-1, +1\}\) y sea
    \((x_{t})_{t \in [0, T]}\) una variable aleatoria que sigue
    la ecuación \eqref{eq:additive_gaussian_noise_model}. Se trata precisamente de un
    modelo de mezcla gaussiana degenerada con distribuciones a priori iguales a
    \(\frac{1}{2}\), medias \(\{-1, +1\}\) y covarianzas ambas iguales
    a \(0\). Para un valor fijo de \(t > 0\), podemos utilizar el cálculo del
    filtro óptimo de ruido de Bayes en \eqref{eq:gmm_bayes_optimal_denoiser}
    para obtener (la demostración se deja como ejercicio)
    \begin{equation}
        \bar{x}^{\star}(t, x_{t}) = \frac{\phi(x_{t} ; +1, t^{2}) - \phi(x_{t}
        ; -1, t^{2})}{\phi(x_{t} ; 1, t^{2}) + \phi(x_{t} ; -1,
        t^{2})} = \tanh\rp{\frac{x_{t}}{t^{2}}}.
    \end{equation}
    Para valores de \(t\) cercanos a \(0\), esta cantidad se aproxima a \(\{-1, +1\}\) para
    casi todas las entradas \(\bar{x}^{\star}(t, x_{t})\). Sin embargo, para
    \(t\) grande, esta cantidad no se encuentra necesariamente ni siquiera aproximadamente
    en el soporte original de \(x\), que, como recordaremos, es \(\{-1,
    +1\}\). En particular, para \(x_{t} \approx 0\) se cumple que
    \(\bar{x}^{\star}(t, x_{t}) \approx 0\), lo cual se encuentra completamente
    entre los dos puntos posibles. Por lo tanto, \(\bar{x}^{\star}\)
    \textit{no dará como resultado un \(x\) «realista»}. Para ser más precisos,
    la distribución de \(\bar{x}^{\star}(t, x_{t})\) es muy
    diferente de la distribución de \(x\).
\end{example}

\subsubsection{Reducción de ruido incremental}
Por lo tanto, para eliminar el ruido de la muestra muy ruidosa \(\vx_{T}\) (donde,
como recordaremos, \(T\) es el tiempo máximo), no podemos utilizar el filtro de eliminación de ruido solo
una vez. En cambio, debemos aplicarlo muchas veces, de forma análoga al descenso por gradiente
con pasos decrecientes, para converger hacia un punto estacionario
\(\hat{\vx}\). Es decir, utilizamos el filtro de eliminación de ruido para pasar de \(\vx_{T}\) a
\(\hat{\vx}_{T - \delta}\), que se aproxima a \(\vx_{T - \delta}\),
luego de \(\hat{\vx}_ {T - \delta}\) a \(\hat{\vx}_{T - 2\delta}\),
y así sucesivamente, desde \(\hat{\vx}_{\delta}\) hasta \(\hat{\vx} =
\hat{\vx}_{0}\). Cada paso de eliminación de ruido hace que la acción del filtro de eliminación de ruido
se parezca más a un paso de gradiente sobre la densidad (logarítmica) original.

De manera más formal, discretizamos uniformemente el intervalo \([0, T]\) en \(L + 1\)
pasos temporales \(0 = t_{0} < t_{1} < \cdots < t_{L} = T\), es decir,
\begin{equation}
    t_{\ell} = \frac{\ell}{L}T, \qquad \ell \in \{0, 1, \dots, L\}.
\end{equation}
A continuación, para cada \(\ell \in [L] = \{1, 2, \dots, L\}\), yendo de
\(\ell = L\) a \(\ell = 1\), ejecutamos la iteración
\begin{align}
    \hat{\vx}_{t_{\ell - 1}}
    &= \Ex[\vx_{t_{\ell - 1}} \mid \vx_{t_{\ell}} = \hat{\vx}_{t_{\ell}}] \\
    &= \Ex[\vx + t_{\ell - 1}\vg \mid \vx_{t_{\ell}} = \hat{\vx}_{t_{\ell}}] \\
    &= \Ex\rs{\vx + t_{\ell - 1}\cdot\frac{\vx_{t_{\ell}} -
    \vx}{t_{\ell}} \given \vx_{t_{\ell}} = \hat{\vx}_{t_{\ell}}} \\
    &= \frac{t_{\ell - 1}}{t_{\ell}}\hat{\vx}_{t_{\ell}} + \bp{1 -
    \frac{t_{\ell - 1}}{t_{\ell}}}\Ex[\vx \mid \vx_{t_{\ell}} =
    \hat{\vx}_{t_{\ell}}] \\
    &= \bp{1 - \frac{1}{\ell}} \hat{\vx}_{t_{\ell}} + \frac{1}{\ell}
    \bar{\vx}^{\star}(t_{\ell}, \hat{\vx}_{t_{\ell}}). \label{eq:iter denoise}
\end{align}

Al inicio de la iteración, cuando \(\ell\) es grande, apenas
confiamos en la salida del filtro de ruido y, en su mayor parte, conservamos la iteración actual.
Esto tiene sentido, ya que el filtro de eliminación de ruido puede presentar una gran varianza (véase
el \Cref{example:denoising_twopoints}). Cuando \(\ell\) es pequeño, el
filtro de ruido se «fija» en los modos de la distribución de los datos, ya que un
paso de la eliminación de ruido es, en esencia, un paso de gradiente sobre la
densidad logarítmica de la distribución verdadera y podemos confiar en que no generará
muestras poco realistas; por lo tanto, la etapa de eliminación de ruido viene determinada por la
salida del filtro de ruido. Cuando \(\ell = 1\), incluso descartamos la
iteración actual y conservamos únicamente la salida del filtro de ruido.

\begin{figure}[t]
    \centering
    \includegraphics[width=\textwidth]{\toplevelprefix/chapters/chapter3/figs/ve_gmm_denoising.png}
    \caption{\small\textbf{Eliminación del ruido en una mezcla de gaussianas
        de rango bajo.} Cada figura visualiza muestras de la distribución real de los datos
        (gris, naranja, rojo) y muestras sometidas al
        proceso de eliminación de ruido \eqref{eq:denoising-iteration-basic} (azul
        claro). En la parte superior izquierda, el proceso acaba de comenzar y el ruido
        es muy intenso. A medida que avanza el proceso, el ruido se desplaza
        cada vez más hacia el soporte de la distribución de datos de rango bajo.
        Por último, en la parte inferior derecha, las muestras generadas se alinean perfectamente
        con el soporte de los datos y se asemejan mucho a
    las muestras extraídas del modelo de mezcla gaussiana de rango bajo.}
    \label{fig:ve_gmm_denoising}
\end{figure}

Visualizamos el proceso de convergencia en \(\R^{3}\) en la
\Cref{fig:ve_gmm_denoising}. Aunque más adelante desarrollaremos algunos resultados rigurosos
sobre la convergencia, ahora ofrecemos un ejemplo en forma cerrada del
proceso iterativo de eliminación de ruido aplicado a una distribución de datos gaussiana.

\begin{example}[Reducción de ruido iterativa para una sola gaussiana]
    \label{ex:single_gaussian_denoising}
    Para comprender de manera intuitiva la iteración de eliminación de ruido anterior,
    la aplicamos para eliminar el ruido de \(\vx_{T}\), definido en
    \eqref{eq:additive_gaussian_noise_model}, donde \(\vx \sim
    \dNorm(\vmu, \vSigma)\). Utilizando \eqref{eq:iter denoise}, calculamos
    \begin{align}
        \hat{\vx}_{t_{\ell - 1}} - \vmu
        &= \bp{1 - \frac{1}{\ell}}\hat{\vx}_{t_{\ell}} +
        \frac{1}{\ell}\bar{\vx}^{\star}(t_{\ell},
        \hat{\vx}_{t_{\ell}}) - \vmu \\
        &= \bp{1 - \frac{1}{\ell}}(\hat{\vx}_{t_{\ell}} - \vmu) +
        \frac{1}{\ell}(\bar{\vx}^{\star}(t_{\ell},
        \hat{\vx}_{t_{\ell}}) - \vmu).
    \end{align}
    A continuación, al sustituir la expresión de \(\bar{\vx}^{\star}\) que
    obtenemos de \eqref{eq:gaussian_denoiser} y realizar algunas
    operaciones tediosas de álgebra matricial, obtenemos
    \begin{equation}
        \vx_{t_{\ell - 1}} - \vmu = \vV\mat{1 - \frac{\ell
            T^{2}}{\lambda_{1}L^{2} + \ell^{2} T^{2}} & & \\ & \ddots &
            \\ & & 1 - \frac{\ell T^{2}}{\lambda_{D}L^{2} +
        \ell^{2}T^{2}}}\vV^{\top}(\vx_{t_{\ell}} - \vmu).
    \end{equation}
    Obsérvese que todas las entradas de esta matriz diagonal están en el intervalo \((0,
    1)\). Si tomamos la norma \(\ell^{2}\) y utilizamos un límite superior sencillo pero estricto, 
    se cumple que
    \begin{equation}\label{eq:denoising_contraction_1step}
        \norm{\vx_{t_{\ell - 1}} - \vmu}_{2} \leq \underbrace{\max_{i
            \in [D]}\bc{1 - \frac{\ell T^{2}}{\lambda_{i}L^{2} +
        \ell^{2}T^{2}}}}_{c_{\ell, L}} \cdot \norm{\vx_{t_{\ell}} - \vmu}_{2}.
    \end{equation}
    A primera vista, esta desigualdad parece implicar que la
    iteración de eliminación de ruido es, conceptualmente, \textit{muy similar} a un
    \textit{mapeo de contracción} como la iteración de potencias
    \eqref{eq:power_iteration} que envía el iterado \(\vx_{t}\) a
    la media \(\vmu\), ya que \(c_{\ell, L}\) están todos en \((0, 1)\)
    para cualquier \(\ell\) y \(L\). Sin embargo, esto es falso. Para ver por qué,
    fijarse en que, al iterar \eqref{eq:denoising_contraction_1step}, se cumple que
    \begin{equation}
        \norm{\vx_{0} - \vmu}_{2} \leq \prod_{\ell = 1}^{L}c_{\ell,
        L} \cdot \norm{\vx_{T} - \vmu}_{2}.
    \end{equation}
    Si cada actualización de eliminación de ruido tuviera propiedades similares a las de un
    mapeo de contracción, entonces \(\prod_{\ell = 1}^{L}c_{\ell, L}\) tendería a
    \(0\) cuando \(L \to \infty\) y la iteración convergería a
    \(\vmu\) extremadamente rápido, lo que haría que el muestreador fuera completamente inútil
    o careciera de sentido. Pero, en realidad, la iteración de eliminación de ruido
    \textit{no es, en realidad, una aplicación de contracción}, ya que \(c_{\ell, L}\)
    varía con ambos índices \(\ell\) y \(L\), de modo que \(\prod_{\ell =
    1}^{L}c_{\ell, L}\) no tiende a \(0\) cuando \(L \to \infty\), como
    se demuestra en la \Cref{fig:not_contraction_mapping} (también es
    posible calcular el límite analíticamente). En realidad, el
    muestreador garantiza que las iteraciones finales se distribuyan a lo largo del
    soporte de la distribución de datos, lo que permite un muestreo real a partir de
    la distribución de datos, algo muy distinto del objetivo de
    la iteración de potencias (por ejemplo).
    \begin{figure}
        \centering
        \includegraphics[width=\textwidth]{\toplevelprefix/chapters/chapter3/figs/not_contraction_mapping.png}
        \caption{\small\textbf{La iteración de reducción de ruido en el \Cref{ex:single_gaussian_denoising} 
                no es un mapeo de 
            contracción.} Gráfico de los productos de los coeficientes de contracción
            \(\prod_{\ell = 1}^{L}c_{\ell, L}\) en el caso en que los
            datos se encuentran en \(D = 1\), el horizonte temporal \(T = 1\) y la
            varianza de los datos \(\lambda_{1} = a\), para distintos valores de
            \(a\). Si la iteración de eliminación de ruido fuera (conceptualmente similar
            a) una aplicación de contracción, estos productos tenderían a \(0\),
        pero, en cambio, convergen hacia valores distintos de cero.}
        \label{fig:not_contraction_mapping}
    \end{figure}
\end{example}

Recordemos que queríamos construir el proceso iterativo de eliminación de ruido para
reducir la entropía de las iteraciones. Aunque lo hemos hecho de forma indirecta,
invirtiendo un proceso que añade entropía, ahora es el momento de
pagar las consecuencias y confirmar que nuestro proceso iterativo de eliminación de ruido reduce la
entropía.

\begin{theorem}[Versión simplifiada del \Cref{thm:conditioning_reduces_entropy}]
    Supongamos que \((\vx_{t})_{t \in [0, T]}\) obedece a
    \eqref{eq:additive_gaussian_noise_model}. Entonces, bajo ciertas
    condiciones técnicas sobre \(\vx\), para todo \(s < t\) con \(s, t
    \in (0, T]\),
    \begin{equation}
        h(\Ex[\vx_{s} \mid \vx_{t}]) < h(\vx_{t}).
    \end{equation}
\end{theorem}
La formulación completa del teorema y la demostración en sí requieren
cierto grado de tecnicismo, por lo que se posponen hasta la
\Cref{sub:denoising_entropy_decreases}.

\subsubsection{Cálculo del filtro de ruido}
Lo último que comentamos aquí es que, en muchas ocasiones,
\textit{no podremos calcular} \(\bar{\vx}^{\star}(t, \cdot)\) para
ningún \(t\), ya que no disponemos de la distribución \(p_{t}\). Pero podemos
intentar \textit{aprender una a partir de los datos}. Recordemos que el filtro de ruido
\(\bar{\vx}^{\star}\) se define en \eqref{eq:denoising_loss} como
la minimización del error cuadrático medio \(\Ex\norm{\bar{\vx}(t, \vx_{t}) -
\vx}_{2}^{2}\). Podemos utilizar este error cuadrático medio como función de pérdida o
función objetivo para entrenar el filtro de reducción de ruido. Por ejemplo, podemos
parametrizar \(\bar{\vx}(t, \cdot)\) mediante una red neuronal, escribiéndola
como \(\bar{\vx}_{\theta}(t, \cdot)\) y optimizar la pérdida en el
espacio de parámetros \(\Theta\):
\begin{equation}
    \min_{\theta \in \Theta}\Ex\norm{\bar{\vx}_{\theta}(t, \vx_{t}) -
    \vx}_{2}^{2}.
\end{equation}
La solución a este problema de optimización, implementada mediante descenso por gradiente
o un algoritmo similar, nos dará un
\(\bar{\vx}_{\theta^{\star}}(t, \cdot)\) que es una buena aproximación
a \(\bar{\vx}^{\star}(t, \cdot)\) (al menos si el entrenamiento funciona)
y que utilizaremos como nuestro filtro de ruido.

¿Cuál sería una buena arquitectura para esta red neuronal
\(\bar{\vx}_{\theta^{\star}}(t, \cdot)\)? Para responder a esta pregunta,
intentaremos primero comprender el caso en que los datos \(\vx\) son
de \textit{bajo rango}, es decir, que tienen un soporte de bajo rango. La forma más canónica
de representar datos de rango bajo es mediante una \textit{distribución gaussiana de rango bajo
}, por lo que comenzamos nuestro análisis suponiendo que \(\vx\)
se extrae (aproximadamente) de una distribución gaussiana de rango bajo:
\begin{equation}\label{eq:x_gaussian_distribution}
    \vx \sim \dNorm(\vzero, \vU\vU^{\top}),
\end{equation}
donde \(\vU \in \O(D, P)\) es una matriz ortogonal alta con \(P \leq
D\) (utilizando la media cero básicamente por conveniencia notacional). El
filtro de ruido óptimo bajo \eqref{eq:x_gaussian_distribution} es (según
el \Cref{example:denoising_gaussian_mixture}):
\begin{equation}\label{eq:optimal denoiser}
    \bar{\vx}^{\star}(t, \vx_{t}) = \vU\vU^{\top}(\vU\vU^{\top} +
    t^{2}\vI)^{-1}\vx_{t}.
\end{equation}
La matriz  \(\vU\vU^{\top} + t^{2}\vI\) es una perturbación de bajo rango
de la matriz de rango completo \(t^{2}\vI\), por lo que se presta a
una simplificación mediante la \textit{identidad de Sherman-Morrison-Woodbury},
es decir, para matrices \(\vA, \vC, \vU, \vV\) tales que \(\vA\) y
\(\vC\) son invertibles,
\begin{equation}\label{eq:sherman_morrison_woodbury}
    (\vA + \vU\vC\vV)^{-1} = \vA^{-1} - \vA^{-1}\vU(\vC^{-1} +
    \vV\vA^{-1}\vU)^{-1}\vV\vA^{-1}.
\end{equation}
Demostramos esta identidad en el
\Cref{exercise:sherman_morrison_woodbury_identity}.  Por ahora, aplicamos
esta identidad con \(\vA = t^{2}\vI\), \(\vU = \vU\), \(\vV =
\vU^{\top}\) y \(\vC = \vI\), obteniendo
\begin{align}
    (\vU\vU^{\top} + t^{2}\vI)^{-1}
    &= \frac{1}{t^{2}}\vI - \frac{1}{t^{4}}\vU\bp{\vI +
    \frac{1}{t^{2}}\vU^{\top}\vU}^{-1}\vU^{\top} \\
    &= \frac{1}{t^{2}}\vI - \frac{1}{t^{4}\bp{1 +
    \frac{1}{t^{2}}}}\vU\vU^{\top} \\
    &= \frac{1}{t^{2}}\bp{\vI - \frac{1}{1 +
    t^{2}}\vU\vU^{\top}}.\label{eq:gaussian_lowrank_inverse_perturbation_simplified}
\end{align}
Esto implica
\begin{align}
    \vU\vU^{\top}(\vU\vU^{\top} + t^{2}\vI)^{-1}
    &= \frac{1}{t^{2}}\vU\vU^{\top}\bp{\vI - \frac{1}{1 +
    t^{2}}\vU\vU^{\top}} \\
    &= \frac{1}{t^{2}}\bp{1 - \frac{1}{1 + t^{2}}}\vU\vU^{\top} \\
    &= \frac{1}{1 +
    t^{2}}\vU\vU^{\top}.\label{eq:gaussian_lowrank_denoiser_simplified}
\end{align}
Nuestro filtro de eliminación de ruido óptimo es una proyección lineal (reescalada) sobre el
subespacio de rango bajo que sustenta la distribución de nuestros datos. Esto es
básicamente lo mismo que el PCA probabilístico del
\Cref{ch:classic-models} (véase ~\Cref{subsec:probabilistic PCA}). Podemos
formalizar esta intuición considerando una clase de
filtros de ruido parametrizados por subespacios:
\begin{equation}\label{eq:optimal nn}
    \bar{\vx}_{\vV}(t, \vx_{t}) = \frac{1}{1 + t^{2}}\vV\vV^{\top}\vx_{t},
\end{equation}
y mostrando que el problema de eliminación de ruido \(\min_{\vV \in \O(D,
    P)}\Ex\norm{\bar{\vx}_{\vV}(t, \vx_{t}) -
\vx}_{2}^{2}\) obtiene la misma solución que el PCA probabilístico
(\Cref{subsec:probabilistic PCA}). De hecho, al sustituir nuestro filtro de ruido por subespacios
en la pérdida de entrenamiento \eqref{eq:denoising_loss}, obtenemos
\begin{equation}
    \min_{\vV \in \O(D, P)}\Ex\norm{\bar{\vx}_{\vV}(t,
    \vx_{t}) - \vx}_{2}^{2} = \Ex\norm*{\frac{1}{1 +
    t^{2}}\vV\vV^{\top}(\vx + t\vg) - \vx}_{2}^{2},
\end{equation}
donde la igualdad se debe a
\eqref{eq:additive_gaussian_noise_model}. Condicionando respecto a \(\vx\)
(integrando así sobre \(\vg\)), calculamos
\begin{align}
    & \Ex_{\vg}\norm*{\frac{1}{1 + t^{2}}\vV\vV^{\top}(\vx + t\vg) -
    \vx}_{2}^{2} \\
    &= \norm*{\frac{1}{1 + t^{2}}\vV\vV^{\top}\vx - \vx}_{2}^{2} -
    \frac{t}{1+t^2}\Ex_{\vg} \ip*{\vx - \frac{1}{1 +
    t^{2}}\vV\vV^{\top}\vx}{\vV\vV^{\top}\vg} \\
    &\quad + \frac{t^2}{(1+t^2)^2} \Ex_{\vg}\norm*{\vV\vV^{\top} \vg}_{2}^{2} \\
    &= \norm*{\frac{1}{1 + t^{2}}\vV\vV^{\top}\vx - \vx}_{2}^{2} +
    \frac{t^2P}{(1+t^2)^2}
\end{align}
donde la segunda igualdad se deduce de \(\vg \sim \dNorm(\vzero,
\vI)\) y del hecho de que \(\vV \in \O(D, P)\), de modo que \(
    \Ex_{\vg}\norm{\vV\vV^\top
    \vg}_{2}^{2} =
\Ex_{\vg}\norm{\vV^\top \vg}_{2}^{2} = P\). Por lo tanto, el problema de eliminación de ruido
es equivalente a
\begin{align}
    \min_{\vV \in \O(D, P)} & \Bigg\{\Ex \norm*{\frac{1}{1 +
        t^{2}}\vV\vV^{\top}\vx - \vx}_{2}^{2} \\
        & \qquad = \Ex\norm{\vx}_{2}^{2} + \bp{\frac{1}{(1 +
            t^{2})^{2}} - \frac{2}{1 +
    t^{2}}}\Ex\norm{\vV^{\top}\vx}_{2}^{2}\Bigg\},
\end{align}
lo cual, a su vez, equivale al problema
\begin{equation}\label{eq:PCA}
    \max_{\vV \in \O(D, P)}\Ex\norm{\vV^{\top}\vx}_{2}^{2}
\end{equation}
lo cual se demostró en la \Cref{sub:pca} que es equivalente al problema
\begin{equation}
    \min_{\vV \in \O(D, P)}\Ex\norm{\vx - \vV\vV^{\top}\vx}_{2}^{2},
\end{equation}
es decir, la formulación probabilística del PCA \eqref{eq:PPCA}.

Por lo tanto, si supusiéramos que nuestros datos proceden de una distribución gaussiana de rango bajo,
nuestro algoritmo de eliminación de ruido consistiría básicamente en aplicar sucesivamente
muchas iteraciones de PCA a datos cada vez menos ruidosos.
Aunque este método parece muy sencillo, en realidad puede generar comportamientos de eliminación de ruido complejos
y relativamente potentes cuando se aplica a datos reales,
tal y como ha demostrado PCANet \parencite{chan2015pcanet}.
No obstante, para obtener un operador de eliminación de ruido más general
que pueda funcionar con distribuciones de datos complejas y multimodales,
consideraremos una clase más amplia de distribuciones de datos: la omnipresente
clase de \textit{modelos de mezcla gaussiana de bajo rango}, cuyo filtro de ruido
calculamos en el \Cref{ex:gaussian_mixture} y el
\Cref{example:denoising_gaussian_mixture}. En el contexto de la
eliminación de ruido, dado que los modelos de mezcla gaussiana pueden aproximar una clase muy amplia
de distribuciones con una precisión arbitraria, en la práctica, la optimización entre
la clase de filtros de ruido para modelos de mezcla gaussiana puede, potencialmente,
proporcionarnos algo cercano al filtro de ruido óptimo para la distribución real de los datos.

Formalmente, supongamos ahora que \(\vx\) se extrae de un modelo de mezcla gaussiana de bajo rango
\(\frac{1}{K}\sum_{k = 1}^{K}\dNorm(\vzero,
\vU_{k}\vU_{k}^{\top})\), donde \(\vU_{k} \in \O(D, P)\) son matrices ortogonales altas
cuyas columnas abarcan los subespacios de dimensión \(P\)
que soportan los \(K\) componentes gaussianos de rango bajo. Escribimos
\begin{equation}\label{eq:MoG1}
    \vx \sim \frac{1}{K}\sum_{k = 1}^{K}\dNorm(\vzero, \vU_{k}\vU_{k}^{\top}).
\end{equation}
Entonces, el filtro de ruido óptimo según
\eqref{eq:additive_gaussian_noise_model} es (según
el \Cref{example:denoising_gaussian_mixture})
\begin{align}
    \bar{\vx}^{\star}(t, \vx_{t})
    &= \sum_{k = 1}^{K}\Bigg\{\frac{\phi(\vx_{t} ; \vzero,
        \vU_{k}\vU_{k}^{\top} + t^{2}\vI)}{\sum_{i = 1}^{K}\phi(\vx_{t} ;
                \vzero, \vU_{i}\vU_{i}^{\top} +
        t^{2}\vI)} \\
        &\qquad \qquad \qquad \qquad
        \cdot\bp{\vU_{k}\vU_{k}^{\top}(\vU_{k}\vU_{k}^{\top} +
    t^{2}\vI)^{-1}\vx_{t}}\Bigg\}.
\end{align}
Al aplicar la misma identidad de Sherman-Morrison-Woodbury que en
\eqref{eq:gaussian_lowrank_inverse_perturbation_simplified}, obtenemos
\begin{equation}
    (\vU_{k}\vU_{k}^{\top} + t^{2}\vI)^{-1} = \frac{1}{t^{2}}\bp{\vI
    - \frac{1}{1 + t^{2}}\vU_{k}\vU_{k}^{\top}}.
\end{equation}
Podemos entonces calcular las probabilidades a posteriori de la siguiente manera. Dado que los
\(\vU_{k}\) son ortogonales, \(\det(\vU_{k}\vU_{k}^{\top}
+ t^{2}\vI)\) es igual para todo \(k\). Por lo tanto
\begin{align}
    &\frac{\phi(\vx_{t} ; \vzero, \vU_{k}\vU_{k}^{\top} + t^{2}\vI)}{\sum_{i
    = 1}^{K}\phi(\vx_{t} ; \vzero, \vU_{i}\vU_{i}^{\top} + t^{2}\vI)} \\
    &= \frac{\exp\rp{-\frac{1}{2}\vx_{t}^{\top}(\vU_{k}\vU_{k}^{\top}
    + t^{2}\vI)^{-1}\vx_{t}}}{\sum_{i =
        1}^{K}\exp\rp{-\frac{1}{2}\vx_{t}^{\top}(\vU_{i}\vU_{i}^{\top} +
    t^{2}\vI)^{-1}\vx_{t}}} \\
    &= \frac{\exp\rp{-\frac{1}{2t^{2}}\vx_{t}^{\top}\bp{\vI -
    \frac{1}{1 + t^{2}}\vU_{k}\vU_{k}^{\top}}\vx_{t}}}{\sum_{i =
        1}^{K}\exp\rp{-\frac{1}{2t^{2}}\vx_{t}^{\top}\bp{\vI - \frac{1}{1
    + t^{2}}\vU_{i}\vU_{i}^{\top}}\vx_{t}}} \\
    &= \frac{\exp\rp{-\frac{1}{2t^{2}}\norm{\vx_{t}}_{2}^{2} +
            \frac{1}{2t^{2}(1 +
    t^{2})}\norm{\vU_{k}^{\top}\vx_{t}}_{2}^{2}}}{\sum_{i =
        1}^{K}\exp\rp{-\frac{1}{2t^{2}}\norm{\vx_{t}}_{2}^{2} +
    \frac{1}{2t^{2}(1 + t^{2})}\norm{\vU_{i}^{\top}\vx_{t}}_{2}^{2}}} \\
    &=
    \frac{\exp\rp{-\frac{1}{2t^{2}}\norm{\vx_{t}}_{2}^{2}}\exp\rp{\frac{1}{2t^{2}(1
                    +
    t^{2})}\norm{\vU_{k}^{\top}\vx_{t}}_{2}^{2}}}{\exp\rp{-\frac{1}{2t^{2}}\norm{\vx_{t}}_{2}^{2}}\sum_{i
        = 1}^{K}\exp\rp{\frac{1}{2t^{2}(1 +
    t^{2})}\norm{\vU_{i}^{\top}\vx_{t}}_{2}^{2}}} \\
    &= \frac{\exp\rp{\frac{1}{2t^{2}(1 +
    t^{2})}\norm{\vU_{k}^{\top}\vx_{t}}_{2}^{2}}}{\sum_{i =
        1}^{K}\exp\rp{\frac{1}{2t^{2}(1 +
    t^{2})}\norm{\vU_{i}^{\top}\vx_{t}}_{2}^{2}}}.
\end{align}
Se trata de una operación softmax ponderada por la proyección de \(\vx_{t}\)
sobre cada subespacio medido por \(\norm{\vU_{i}^{\top}\vx_{t}}_{2}\)
(atenuada por la temperatura \(2t^{2}(1 + t^{2})\)). Mientras tanto, los
filtros de ruido componentes se pueden escribir de la misma manera que
\eqref{eq:gaussian_lowrank_denoiser_simplified} para obtener
\begin{equation}
    \vU_{k}\vU_{k}^{\top}(\vU_{k}\vU_{k}^{\top} +
    t^{2}\vI)^{-1}\vx_{t} = \frac{1}{1 + t^{2}}\vU_{k}\vU_{k}^{\top}\vx_{t}.
\end{equation}
Si juntamos todo esto, tenemos
\begin{equation}\label{eq:gmm_lowrank_denoiser}
    \bar{\vx}^{\star}(t, \vx_{t}) = \frac{1}{1 + t^{2}}\sum_{k =
    1}^{K}\frac{\exp\rp{\frac{1}{2t^{2}(1 +
    t^{2})}\norm{\vU_{k}^{\top}\vx_{t}}_{2}^{2}}}{\sum_{i =
        1}^{K}\exp\rp{\frac{1}{2t^{2}(1 +
    t^{2})}\norm{\vU_{i}^{\top}\vx_{t}}_{2}^{2}}}\vU_{k}\vU_{k}^{\top}\vx_{t},
\end{equation}
es decir, una proyección de \(\vx_{t}\) sobre cada uno de los \(K\) subespacios,
ponderada por una función softmax de una función cuadrática de \(\vx_{t}\). ¡Esta
forma funcional es similar a un \textit{mecanismo de atención} en una 
arquitectura de transformador! Como veremos en el
\Cref{ch:deep-networks}, esto no es en absoluto una coincidencia; el profundo
vínculo entre la eliminación de ruido y la compresión con pérdida (que se tratará en la
\Cref{sec:lossy_compression}) hace que los filtros de ruido basados en transformadores sean tan
eficaces en la práctica.

En general, la arquitectura del filtro de ruido aprendido corresponde a un
esquema de codificación (paramétrico implícito) de los datos dados, ya que puede
utilizarse para eliminar el ruido o proyectar sobre la distribución de los datos. Entrenar un
filtro de ruido equivale a encontrar un mejor esquema de codificación, lo que
implementa parcialmente el enfoque A2 (es decir, buscar una
distribución mediante la búsqueda de esquemas con mejores tasas de codificación) al
final de la \Cref{sub:min_entropy}. A continuación, analizamos cómo
implementar el otro enfoque.

\subsection{Muestreo de una distribución mediante
eliminación iterativa de ruido}\label{sub:sampling_denoising}

Recordemos que, al final de la \Cref{sub:min_entropy}, analizamos dos
enfoques para buscar una distribución con estructura de baja dimensión.
El primer enfoque (A1) comienza con una distribución normal de alta
entropía y reduce gradualmente su entropía hasta alcanzar la distribución
de los datos. A este procedimiento lo llamamos \textit{muestreo}, ya que
generamos nuevas muestras. Ha llegado el momento de analizar cómo poner esto en práctica
con el conjunto de herramientas que hemos creado.

Sabemos cómo eliminar el ruido de muestras muy ruidosas \(\vx_{T}\) para obtener
aproximaciones \(\hat{\vx}\) cuyas distribuciones se asemejan a la de la
variable aleatoria objetivo \(\vx\). Sin embargo, el procedimiento de muestreo nos obliga
a partir de una distribución modelo que \textit{no} reciba influencia
de la distribución de \(\vx\) y a utilizar el filtro de ruido para guiar las
iteraciones hacia la distribución de \(\vx\). ¿Cómo podemos hacerlo? Una
forma de hacerlo se explica a continuación:
\begin{equation}\label{eq:ve_converge_to_large_gaussian}
    \frac{\vx_{T}}{T} = \frac{\vx + T\vg}{T} = \frac{\vx}{T} + \vg
    \xrightarrow{T\to +\infty} \vg \sim \dNorm(\vzero, \vI).
\end{equation}
Por lo tanto, \(\vx_{T} \approx \dNorm(\vzero, T^{2}\vI)\). Esta
aproximación es bastante buena para casi todas las distribuciones prácticas,
y se muestra en la \Cref{fig:xT_vs_noise}.
\begin{figure}
    \centering
    \includegraphics[width=0.75\textwidth]{\toplevelprefix/chapters/chapter3/figs/xT_vs_noise.png}
    \caption{\small\textbf{Visualización de \(\vx_{T}\) vs
        \(\dNorm(\vzero, T^{2}\vI)\).} \textit{Izquierda:} Gráfico de los datos del modelo de mezcla gaussiana
        \(\vx\). \textit{Derecha:} Gráfico de \(\vx\),
        así como de \(\vx_{T}\) y una muestra independiente de \(\dNorm(\vzero,
        T^{2}\vI)\), para \(T = 10\). En el gráfico de la derecha, \(\vx\) se
        representa con los mismos colores que en el de la izquierda; sin embargo, las muestras de
        \(\vx_{T}\) y \(\dNorm(\vzero, T^{2}\vI)\) son, en promedio, mucho
        más grandes que las muestras de \(\vx\), por lo que este último parece
        mucho más pequeño debido a la escala. A pesar de esta gran ampliación,
        observamos claramente las similitudes en las distribuciones de
    \(\vx_{T}\) y \(\dNorm(\vzero, T^{2}\vI)\).} 
    \label{fig:xT_vs_noise}
\end{figure}

Así pues, al discretizar \([0, T]\) en \(0 = t_{0} < t_{1} < \cdots < t_{L}
= T\) de manera uniforme utilizando \(t_{\ell} = T\ell / L\) (como en la sección anterior),
una forma posible de obtener muestras a partir de ruido puro es:
\begin{itemize}
    \item Muestrear \(\hat{\vx}_{T} \sim \dNorm(\vzero, T^{2}\vI)\)
        (independientemente de todo lo demás).
    \item Ejecutar la iteración de eliminación de ruido tal y como se indica en la
        \Cref{sub:intro_diffusion_denoising}, es decir, 
        \begin{equation}\label{eq:denoising-iteration-basic}
            \hat{\vx}_{t_{\ell - 1}} = \bp{1 -
            \frac{1}{\ell}}\hat{\vx}_{t_{\ell}} +
            \frac{1}{\ell}\bar{\vx}^{\star}(t_{\ell}, \hat{\vx}_{t_{\ell}}).
        \end{equation}
    \item Output \(\hat{\vx} = \hat{\vx}_{0}\).
\end{itemize}
Conceptualmente, esto es todo lo que hay detrás de los \textit{modelos de difusión},
que transforman el ruido en muestras de datos de acuerdo con el primer
requisito. Sin embargo, aún quedan algunos pasos por dar antes de obtener modelos que
puedan realmente tomar muestras de distribuciones de datos reales, como imágenes, bajo
restricciones prácticas de recursos. En lo que sigue, presentamos y
justificamos varios de esos pasos.

\paragraph{Paso 1: Diferentes discretizaciones.} El primer paso se basa
en la observación de que \textit{no necesitamos realizar tantas
iteraciones de eliminación de ruido} para valores grandes de \(t\). La \Cref{fig:ve_gmm_denoising}
muestra que las primeras \(200\) o \(300\) iteraciones de las \(500\)
iteraciones del proceso de muestreo se dedican a reducir el ruido
hacia la distribución de los datos en su conjunto, antes de que las iteraciones restantes
empujen las muestras hacia un subespacio. Dado un número fijo
de iteraciones \(L\), esto sugiere que deberíamos asignar más
pasos temporales \(t_{\ell}\) cerca de \(t = 0\) que cerca de \(t = T\). Durante
el muestreo (y el entrenamiento), podemos, por lo tanto, utilizar una
discretización alternativa de \([0, T]\) en \(0 \leq t_{0} < t_{1} < \cdots <
t_{L} \leq T\), como por ejemplo una \textit{discretización exponencial}: 
\begin{equation}\label{eq:denoising_exponential_discretization}
    t_{\ell} = C_{1}(e^{C_{2}\ell} - 1), \qquad \forall \ell \in \{0,
    1, \dots, L\}
\end{equation}
donde \(C_{1}, C_{2} > 0\) son constantes que pueden ajustarse para
obtener un rendimiento óptimo en la práctica; los análisis teóricos suelen especificar
también dichas constantes óptimas. La iteración de eliminación de ruido/muestreo se convierte entonces
en
\begin{equation}\label{eq:denoising_iteration}
    \hat{\vx}_{t_{\ell - 1}} \doteq \frac{t_{\ell -
    1}}{t_{\ell}}\hat{\vx}_{t_{\ell}} + \bp{1 - \frac{t_{\ell -
    1}}{t_{\ell}}}\bar{\vx}^{\star}(t_{\ell}, \hat{\vx}_{t_{\ell}}),
\end{equation}
con \(\hat{\vx}_{t_{L}} \sim \dNorm(\vzero, t_{L}^{2}\vI)\).

\begin{figure}[t]
    \centering
    \includegraphics[width=\textwidth]{\toplevelprefix/chapters/chapter3/figs/vp_gmm_denoising.png}
    \caption{\small\textbf{Eliminación del ruido en una mezcla de distribuciones gaussianas mediante el
        proceso de difusión VP.} Utilizamos la misma configuración de la figura y la misma distribución de datos
        que en la \Cref{fig:ve_gmm_denoising}. Nótese que, en comparación
        con la \Cref{fig:ve_gmm_denoising}, la distribución del ruido está mucho
    más concentrada alrededor del origen.}
    \label{fig:vp_gmm_denoising}
\end{figure}

\paragraph{Paso 2: Modelos de ruido diferentes.} El segundo paso consiste en
analizar modelos ligeramente diferentes a los de
\eqref{eq:additive_gaussian_noise_model}. La razón fundamental para
esto es la siguiente. En la práctica, la distribución del ruido
\(\dNorm(\vzero, t_{L}^{2}\vI)\) se convierte en una estimación cada vez
menos precisa de la covarianza real en dimensiones elevadas; es decir,
\eqref{eq:ve_converge_to_large_gaussian} se convierte en una aproximación cada vez
menos precisa, especialmente con datos anisotrópicos de alta dimensión.
El aumento de la distancia entre \(\dNorm(\vzero, t_{L}^{2}\vI)\)
y la distribución real de \(\vx_{t_{L}}\) puede hacer que el filtro de ruido
funcione peor en tales circunstancias. En teoría,
\(\vx_{t_{L}}\) nunca converge a ninguna distribución a medida que \(t_{L}\)
aumenta, por lo que esta configuración es difícil de analizar de principio a fin. En este
caso, nuestra solución consiste en \textit{añadir ruido y reducir
    la contribución de \(\vx\) al mismo tiempo, de modo que \(\vx_{T}\) converja cuando \(T
\to \infty\)}. La tasa de ruido añadido se denota por \(\sigma \colon [0,
T] \to \R_{\geq 0}\) y la tasa de contracción se denota por \(\alpha
\colon [0, T] \to \R_{\geq 0}\), de tal manera que \(\sigma\) es
\textit{creciente} y \(\alpha\) es (no estrictamente) \textit{decreciente}, y
\begin{equation}\label{eq:gen_additive_gaussian_noise_model}
    \vx_{t} \doteq \alpha_{t}\vx + \sigma_{t}\vg, \qquad \forall t \in [0, T].
\end{equation}
En la configuración anterior, \(\alpha_{t} = 1\) y \(\sigma_{t} = t\) y
esto se conoce como el \textit{proceso de varianza explosiva (VE)}. Una
opción habitual que reduce la contribución de \(\vx\), tal y como
describimos inicialmente, es aquella en la que \(T = 1\) (de modo que \(t \in [0, 1]\)),
\(\alpha_{t} = \sqrt{1 - t^{2}}\) y \(\sigma_{t} = t\); este es el
\textit{proceso de preservación de la varianza (VP)}. Nótese que, en el proceso VP,
\(\vx_{1} \sim \dNorm(\vzero, \vI)\) exactamente, por lo que basta con
tomar muestras de esta distribución estándar y aplicar el proceso de eliminación 
de ruido de forma iterativa. Como
resultado, el proceso VP es mucho más fácil de analizar teóricamente y
más estable empíricamente.\footnote{¿Por qué utilizar toda la configuración de \(\alpha,
    \sigma\)? Como veremos en el
    \Cref{exercise:implement_denoising_processes}, esta encapsula y
    unifica muchos de los procesos propuestos, incluido el recientemente popular
    proceso denominado \textit{flow matching}. A pesar de ello, los procesos VE y VP
    siguen siendo los más populares tanto desde el punto de vista empírico como teórico
(hasta ahora), por lo que los analizaremos en esta sección.}

Con esta configuración más general, la fórmula de Tweedie \eqref{eq:tweedie} se convierte en
\begin{equation}\label{eq:gen_tweedie}
    \Ex[\vx \mid \vx_{t}] = \frac{1}{\alpha_{t}}\bp{\vx_{t} +
    \sigma_{t}^{2}\nabla \log p_{t}(\vx_{t})}.
\end{equation}
La iteración de eliminación de ruido \eqref{eq:denoising_iteration} se convierte en
\begin{equation}\label{eq:gen_denoising_iteration}
    \hat{\vx}_{t_{\ell - 1}} = \frac{\sigma_{t_{\ell -
    1}}}{\sigma_{t_{\ell}}}\hat{\vx}_{t_{\ell}} + \bp{\alpha_{t_{\ell
        - 1}} - \frac{\sigma_{t_{\ell -
    1}}}{\sigma_{t_{\ell}}}\alpha_{t_{\ell}}}\bar{\vx}^{\star}(t_{\ell},
    \hat{\vx}_{t_{\ell}}).
\end{equation}

\begin{example}[El modelo de mezcla gaussiana]
    Para un modelo de mezcla gaussiana (véase el \Cref{ex:gaussian_mixture}),
    su filtro de eliminación de ruido \eqref{eq:gmm_bayes_optimal_denoiser} se convierte en
    \begin{align}\label{eq:gen_gmm_bayes_optimal_denoiser}
        \bar{\vx}^{\star}(t, \vx_{t})
        &= \sum_{k = 1}^{K}\Bigg\{\frac{\pi_{k}\phi(\vx_{t}
                    ; \alpha_{t}\vmu_{k}, \alpha_{t}^{2}\vSigma_{k}
            + \sigma_{t}^{2}\vI)}{\sum_{i = 1}^{K}\pi_{i}\phi(\vx_{t} ;
                    \alpha_{t}\vmu_{i}, \alpha_{t}^{2}\vSigma_{i} +
            \sigma_{t}^{2}\vI)} \\
            &\qquad \qquad \qquad \cdot\bp{\vmu_{k} +
                \alpha_{t}\vSigma_{k}(\alpha_{t}^{2}\vSigma_{k} +
        \sigma_{t}^{2}\vI)^{-1}(\vx_{t} - \alpha_{t}\vmu_{k})}\Bigg\}. \nonumber
    \end{align}
    La \Cref{fig:vp_gmm_denoising} muestra iteraciones del procedimiento
    de muestreo.
\end{example}

Cabe señalar que la iteración de eliminación de ruido \eqref{eq:gen_denoising_iteration}
da lugar a un algoritmo de muestreo denominado DDIM («Denoising Diffusion
Implicit Model») \cite{song2020denoising}, y es uno de los
algoritmos de muestreo más populares utilizados actualmente en los modelos de difusión. Lo
resumimos en el \Cref{alg:iterative_denoising}.

\begin{algorithm}
    \caption{Muestreo mediante filtro de ruido.}
    \label{alg:iterative_denoising}
    \begin{algorithmic}[1]
        \Require{Una lista ordenada de intervalos de tiempo \(0 \leq t_{0} < \cdots
        < t_{L} \leq T\) que se utilizarán para el muestreo.}
        \Require{Un filtro de eliminación de ruido \(\bar{\vx} \colon \{t_{\ell}\}_{\ell =
        1}^{L} \times \R^{D} \to \R^{D}\).}
        \Require{Funciones de escala y nivel de ruido \(\alpha, \sigma
        \colon \{t_{\ell}\}_{\ell = 0}^{L} \to \R_{\geq 0}\).}
        \Ensure{Una muestra \(\hat{\vx}\), aproximadamente de la
        distribución de \(\vx\).}
        \Function{DDIMSampler}{$\bar{\vx}, (t_{\ell})_{\ell = 0}^{L}$}
        \State{Inicializar \(\hat{\vx}_{t_{L}} \sim\) la distribución aproximada
        de \(\vx_{t_{L}}\)}
        \Statex{\quad\ \texttt{\# VP \(\implies
        \dNorm(\vzero, \vI)\), VE \(\implies \dNorm(\vzero, t_{L}^{2}\vI)\).}}
        \Statex{}
        \For{\(\ell = L, L - 1, \dots, 1\)}
        \State{Compute
            \begin{equation*}
                \hat{\vx}_{t_{\ell - 1}} \doteq \frac{\sigma_{t_{\ell
                - 1}}}{\sigma_{t_{\ell}}}\hat{\vx}_{t_{\ell}} +
                \bp{\alpha_{t_{\ell - 1}} - \frac{\sigma_{t_{\ell -
                1}}}{\sigma_{t_{\ell}}}\alpha_{t_{\ell}}}\bar{\vx}(t_{\ell},
                \hat{\vx}_{t_{\ell}})
            \end{equation*}
        }
        \EndFor
        \State{\Return{\(\hat{\vx}_{t_{0}}\)}}
        \EndFunction
    \end{algorithmic}
\end{algorithm}

\paragraph{Paso 3: Optimización de los procesos de entrenamiento.} Si usáramos el
procedimiento descrito en la \Cref{sub:intro_diffusion_denoising} para entrenar un
filtro de ruido independiente \(\bar{\vx}(t, \cdot)\) para cada instante \(t\) que se
utilizará en el algoritmo de muestreo, \textit{¡tendríamos que entrenar \(L\)
filtros de ruido independientes!} Esto es muy ineficiente: lo habitual es
que tengamos que entrenar \(L\) redes neuronales distintas, lo que requiere \(L\)
veces más tiempo de entrenamiento y memoria de almacenamiento y luego quedarnos atados a
utilizar esos intervalos de tiempo para el muestreo de por siempre. En cambio, podemos
\textit{entrenar una sola red neuronal} para eliminar el ruido en todos los instantes
\(t\), tomando como entrada las variables continuas \(\vx_{t}\) y \(t\)
(en lugar de solo \(\vx_{t}\) como antes). Desde el punto de vista mecánico, nuestra pérdida de entrenamiento
se promedia a lo largo de \(t\), es decir, resuelve el siguiente problema:
\begin{equation}\label{eq:training loss}
    \min_{\theta}\Ex\norm{\bar{\vx}_{\theta}(\tau,
    \vx_{\tau}) - \vx}_{2}^{2},
\end{equation}
donde \(\tau \sim \dUnif([0, T])\), es decir, se extrae uniformemente del intervalo \([0, T]\).
Al igual que en el Paso 1, donde utilizamos más intervalos de tiempo más cercanos a \(t = 0\)
para garantizar un mejor proceso de muestreo, es posible que queramos asegurarnos de que el
filtro de ruido sea de mayor calidad cuanto más cerca de \(t = 0\) y por lo tanto
\textit{ponderar la pérdida} de modo que \(t\) cerca de \(0\) tenga mayor
ponderación.\footnote{En la práctica, la distribución de la \(t\) muestreada también puede
    ser no uniforme, con el fin de garantizar (junto con el
    esquema de ponderación) que la estimación de Monte Carlo de la pérdida tenga
una varianza reducida, lo que mejora considerablemente la estabilidad del entrenamiento.}
Si \(w_{t}\) es el peso en el momento \(t\), la pérdida ponderada tendría la forma
\begin{equation}\label{eq:parametric_denoising_loss}
    \min_{\theta}\Ex[w_{\tau}\norm{\bar{\vx}_{\theta}(\tau,
    \vx_{\tau}) - \vx}_{2}^{2}].
\end{equation}
Una opción razonable para el peso en la práctica es \(w_{t} =
\alpha_{t}^{2}/\sigma_{t}^{2}\). La razón exacta se explicará en el
siguiente párrafo, pero, en general, sirve para dar mayor peso a las pérdidas
correspondientes a \(t\) cerca de \(0\), sin dejar de mantener una estabilidad numérica
razonable. Además, por supuesto, en la práctica no podemos calcular la
esperanza, por lo que utilizamos la media de Monte Carlo, que es la más sencilla,
para estimarla. La serie de cambios que se han realizado aquí
ofrece varias ventajas conceptuales y computacionales: no necesitamos
entrenar múltiples filtros de ruido, podemos entrenar con un solo conjunto de intervalos de tiempo y
realizar el muestreo utilizando un subconjunto (o incluso otros distintos), etc. El proceso completo se
analiza en el \Cref{alg:learning_denoiser}. Sin duda, este proceso
funciona en la práctica; comentamos algunos resultados en las
\Cref{sec:auto-encode-img-gen,sec:cond_gen,sec:object-generation,sec:shape-and-pose,sec:egoallo}.

\begin{algorithm}
    \caption{Aprender un filtro de ruido a partir de datos.}
    \begin{algorithmic}[1]
        \Require{Conjunto de datos \(\cD \subseteq \R^{D}\).}
        \Require{Una lista ordenada de intervalos de tiempo \(0 \leq t_{0} < \cdots
        < t_{L} \leq T\) que se utilizarán para el muestreo.}
        \Require{Una función de ponderación \(w \colon \{t_{\ell}\}_{\ell =
        1}^{L} \to \R_{\geq 0}\).}
        \Require{Funciones de escala y nivel de ruido \(\alpha, \sigma
        \colon \{t_{\ell}\}_{\ell = 0}^{L} \to \R_{\geq 0}\).}
        \Require{Un espacio de parámetros \(\Theta\) y una arquitectura
        de filtro de ruido \(\bar{\vx}_{\theta}\).}
        \Require{Un algoritmo de optimización para los parámetros.}
        \Require{El número de iteraciones de optimización \(M\).}
        \Require{El número de simulaciones de Monte Carlo \(N\) por paso de entrenamiento.}
        \Ensure{Un filtro de eliminación de ruido entrenado \(\bar{\vx}_{\theta^{\star}}\).}
        \Function{TrainDenoiser}{$\cD, \Theta$}
        \State{Inicializar \(\theta^{(1)} \in \Theta\)}
        \For{\(m \in [M]\)}
        \For{\(i \in [N]\)}
        \Statex{\qquad\qquad\ \texttt{\# Extraer una muestra del
        conjunto de datos.}}
        \State{\(\vx^{(m, i)} \sim \cD\)}
        \Statex{}

        \Statex{\qquad\qquad\ \texttt{\# Muestrear un intervalo de tiempo.}}
        \State{\(t^{(m, i)} \simiid \dUnif(\{t_{\ell}\}_{\ell =
        1}^{L})\)}
        \Statex{}

        \Statex{\qquad\qquad\ \texttt{\# Calcular ponderaciones y escalas.}}
        \State{\((\alpha^{(m, i)}, \sigma^{(m, i)}, w^{(m, i)}) =
        (\alpha_{t^{(m, i)}}, \sigma_{t^{(m, i)}}, w_{t^{(m, i)}})\)}
        \Statex{}

        \Statex{\qquad\qquad\ \texttt{\# Muestrear un vector de ruido.}}
        \State{\(\vg^{(m, i)} \simiid \dNorm(\vzero, \vI)\)}
        \Statex{}

        \Statex{\qquad\qquad\ \texttt{\# Calcular la muestra ruidosa.}}
        \State{\(\vx_{t}^{(m, i)} \doteq \alpha^{(m, i)}\vx^{(m, i)}
        + \sigma^{(m, i)}\vg^{(m, i)}\)}
        \EndFor

        \Statex{\qquad\ \texttt{\# Calcular la pérdida de entrenamiento.}}
        \State{\(\hat{\cL}^{(m)} \doteq \displaystyle
                \frac{1}{N}\sum_{i = 1}^{N}w^{(m, i)}\norm{\vx^{(m, i)} -
                    \bar{\vx}_{\theta^{(m)}}(t^{(m, i)}, \vx_{t}^{(m,
        i)})}_{2}^{2}\)}
        \Statex{}
        \Statex{\qquad\ \texttt{\# Actualizar los parámetros del modelo.}}
        \State{\(\theta^{(m + 1)} \doteq
                \texttt{OptimUpdate}^{(m)}(\theta^{(m)},
        \nabla_{\theta^{(m)}}\hat{\cL}^{(m)})\)}
        \EndFor
        \Statex{\quad\ \texttt{\# Devolver el filtro de ruido entrenado.}}
        \State{\Return{\(\bar{\vx}_{\theta^{(M + 1)}}\)}}
        \EndFunction
    \end{algorithmic}
    \label{alg:learning_denoiser}
\end{algorithm}

\paragraph{Paso 4: Cambio del objetivo de estimación.}
Es habitual reorientar todo el proceso de eliminación de ruido
en torno a los \textit{predictores de ruido}, es decir, las estimaciones de \(\Ex[\vg \mid
\vx_{t}]\). En la práctica, los predictores de ruido son ligeramente más fáciles de
entrenar porque su salida es (casi) siempre de un tamaño comparable al de una
variable aleatoria gaussiana, por lo que el entrenamiento es más estable numéricamente, pero
pueden requerir modelos más grandes, ya que predicen una
variable objetivo de dimensión completa en lugar de datos potencialmente de baja dimensión
\citep{li2025back}. Por \eqref{eq:gen_additive_gaussian_noise_model} tenemos
\begin{equation}
    \vx_{t} = \alpha_{t}\Ex[\vx \mid \vx_{t}] + \sigma_{t}\Ex[\vg
    \mid \vx_{t}] \implies \Ex[\vg \mid \vx_{t}] =
    \frac{1}{\sigma_{t}}\bp{\vx_{t} - \alpha_{t}\Ex[\vx \mid \vx_{t}]},
\end{equation}
por lo que cualquier predictor de \(\vx\) puede transformarse en un predictor
de \(\vg\) utilizando la relación anterior, es decir,
\begin{equation}
    \bar{\vg}(t, \vx_{t}) = \frac{1}{\sigma_{t}}\vx_{t} -
    \frac{\alpha_{t}}{\sigma_{t}}\bar{\vx}(t, \vx_{t}),
\end{equation}
y viceversa. Por lo tanto, una buena red para estimar \(\bar{\vg}\) es
lo mismo que una buena red para estimar \(\bar{\vx}\) \textit{más
una conexión residual} (como se ve, por ejemplo, en los transformadores). Sus pérdidas
son también las mismas que las del filtro de ruido, salvo por el factor
\(\alpha_{t}/\sigma_{t}\), es decir,
\begin{equation}
    \Ex[w_{\tau}\norm{\vg - \bar{\vg}(\tau,
    \vx_{\tau})}_{2}^{2}] =
    \Ex\rs{w_{\tau}\frac{\alpha_{\tau}^{2}}{\sigma_{\tau}^{2}}\norm{\vx
    - \bar{\vx}(\tau, \vx_{\tau})}_{2}^{2}},
\end{equation}
recordando que \(\tau \sim \dUnif([0, T])\). Se han propuesto otros objetivos
para diferentes tareas, por ejemplo, \(\Ex[\vv_{t} \mid
\vx_{t}]\) donde \(\vv_{t} = \odv*{\vx_{t}}{t}\) (denominado
predicción \(v\) o predicción de velocidad), etc. Esto resulta ser
equivalente a la eliminación de ruido y la predicción de ruido, pero también resulta fundamentalmente
útil para generalizaciones de la difusión, como el flow matching, que
analizaremos en la \Cref{rem:flow_matching}.

\subsubsection{Tasa de error de muestreo del proceso de eliminación de ruido}
Hemos introducido muchos cambios en nuestro proceso platónico original de generación y eliminación de ruido.
Para asegurarnos de que el nuevo proceso sigue funcionando en
la práctica, podemos calcular ejemplos numéricos (como
en la \Cref{fig:vp_gmm_denoising}). Para asegurarnos de que es
teóricamente sólido, podemos demostrar una \textit{cota de la tasa de error}
para el algoritmo de muestreo, lo que demuestra que la tasa de error es pequeña.
A continuación, presentamos un valor de este tipo extraído de la literatura, que muestra que la
distribución de los resultados del muestreador converge, en la denominada
\textit{distancia de variación total (TV)}, hacia la distribución verdadera. La
distancia TV entre dos variables aleatorias \(\vx\) y
\(\vy\) se define como:\footnote{El supremo aquí se aplica únicamente a
    conjuntos \(A\) \textit{medibles}, no a todos los conjuntos sin excepción. Al igual que con todas las
    demás advertencias relacionadas con la teoría de la medida que aparecen en este libro, no dudes en
ignorarlo a efectos prácticos.}
\begin{equation}
    \TV(\vx, \vy) \doteq \sup_{A \subseteq \R^{d}}\abs*{\Pr[\vx \in
    A] - \Pr[\vy \in A]}.
\end{equation}
Si \(\vx\) y \(\vy\) están muy cerca (uniformemente), entonces el supremo
será pequeño. Así pues, la distancia TV mide la proximidad entre
variables aleatorias. (Se trata, efectivamente, de una métrica, como sugiere su nombre; la demostración es
un ejercicio.)

\begin{theorem}[\citep{li2024d} Teorema 1,
    simplificado]\label{thm:diffusion_sampler_convergence}
    Supongamos que \(\Ex\norm{\vx}_{2} < \infty\). Si a \(\vx\) se
    le elimina el ruido según el proceso VP con una discretización
    exponencial\footnote{La definición exacta es bastante extensa
        en nuestra notación y solo se define hasta diversas
        constantes absolutas, por lo que la omitimos aquí por brevedad. Se encuentra en el
    artículo original \citep{li2024d}.} como en
    \eqref{eq:denoising_exponential_discretization}, la salida
    \(\hat{\vx}\) del \Cref{alg:iterative_denoising} satisface la
    cota de variación total
    \begin{align}\label{eq:diffusion_sampling_error}
        \TV(\vx, \hat{\vx})
        &= \tilde{\cO}\Bigg(\underbrace{\frac{D}{L}}_{\text{discretization
            error}} \\
            & \qquad \qquad + \underbrace{\sqrt{\frac{1}{L}\sum_{\ell =
                    1}^{L}\frac{\alpha_{t_{\ell}}^{2}}{\sigma_{t_{\ell}}^{4}}\Ex\norm{\bar{\vx}^{\star}(t_{\ell},
                        \vx_{t_{\ell}}) - \bar{\vx}(t_{\ell},
            \vx_{t_{\ell}})}_{2}^{2}}}_{\text{average excess error of
        the denoiser}}\Bigg)\nonumber
    \end{align}
    donde \(\bar{\vx}^{\star}\) es el filtro de ruido óptimo de Bayes para
    \(\vx\) y \(\tilde{\cO}\) es una variante de la notación «big-\(\cO\)»
    que ignora los factores logarítmicos en \(L\).
\end{theorem}
La técnica de demostración de alto nivel consiste, como se ha comentado anteriormente, en
limitar el error en cada paso, distinguir las fuentes de error (entre
el error de discretización y el del filtro de ruido) y asegurarse cuidadosamente de que
los errores no se acumulen en exceso (o incluso se anulen entre sí).

Tenga en cuenta que si \(L \to \infty\) y aprendemos correctamente el filtro de ruido óptimo de Bayes
\(\bar{\vx} = \bar{\vx}^{\star}\) (de modo que el error excedente es
\(0\)), entonces el proceso de muestreo en el
\Cref{alg:iterative_denoising} produce una \textit{inversa perfecta (en
distribución)} del proceso de ruido, ya que la tasa de error en el
\Cref{thm:diffusion_sampler_convergence} tiende a
\(0\),\footnote{Se obtienen resultados similares para los procesos VE; los
    conjuntos de herramientas teóricas utilizados para las demostraciones difieren debido a la estructura de los
procesos estocásticos subyacentes.} tal y como se ha argumentado de manera heurística anteriormente.

\begin{remark}
    ¿Qué pasa si los datos son de baja dimensión, es decir, si se encuentran en un
    \textit{subespacio de rango bajo} del espacio de alta dimensión
    \(\R^{D}\)? Si la distribución de los datos tiene soporte compacto---por ejemplo,
    si los datos se normalizan al hipercubo unitario, lo cual suele
    garantizarse como paso previo al procesamiento de datos reales, como
    las imágenes---es posible obtener mejores resultados. Concretamente, los autores de
    \cite{li2024d} también definen una medida de la \textit{dimensión intrínseca
    aproximada} \(\kappa\) utilizando la asintótica del denominado
    número de cobertura, que es extremadamente similar en su intuición (si no
    en su implementación) a la función de tasa-distorsión que se presenta en
    el siguiente capítulo. A continuación, muestran que, al aplicar una pequeña
    modificación concreta al muestreador DDIM del
    \Cref{alg:iterative_denoising} (es decir, perturbando ligeramente los
    coeficientes de actualización), el error de discretización pasa a ser
    \(\tilde{\cO}(\kappa/L)\) en lugar de \(\tilde{\cO}(D/L)\) como en el
    \Cref{thm:diffusion_sampler_convergence}. Por lo tanto, al utilizar este
    algoritmo modificado, \(L\) no tiene por qué ser demasiado grande, incluso cuando
    \(D\) alcanza valores de miles o millones, ya que los datos reales tienen
    una estructura de baja dimensión. Sin embargo, en la práctica utilizamos el muestreador DDIM,
    por lo que \(L\) debería tener una dependencia moderada de \(D\)
    para lograr tasas de error consistentes. La elección concreta de \(L\)
    supone un equilibrio entre la complejidad computacional (por ejemplo, el tiempo de ejecución o
    el consumo de memoria) del muestreo y la complejidad estadística de
    aprender un filtro de ruido para estructuras de baja dimensión. El valor de
    \(L\) suele ser diferente en la fase de entrenamiento (donde un \(L\) mayor
        permite una mejor cobertura del intervalo \([0, T]\), lo que ayuda
    a la red a aprender una relación que se generaliza sobre \(t\))
    y en la fase de muestreo (donde un \(L\) menor implica un muestreo más eficiente).
    Incluso se pueden seleccionar los intervalos de tiempo de forma adaptativa en el momento del muestreo
    para optimizar este equilibrio \cite{bao2022analytic}.
\end{remark}

\begin{remark}
    En otros trabajos, el proceso inverso se define como un desplazamiento hacia atrás
    en el índice temporal \(t\) mediante una ecuación de diferencias explícita, o
    una ecuación diferencial en el límite \(L \to \infty\), o hacia adelante
    en el tiempo utilizando la transformación \(\vy_{t} = \vx_{T - t}\), de tal manera
    que si \(t\) aumenta, entonces \(\vy_{t}\) se acerca a
    \(\vx_{0}\). Cada artículo utiliza su propia notación y a veces
    puede resultar confuso comprender las relaciones entre
    los resultados de los distintos artículos. En este trabajo nos esforzamos por mantener
    la coherencia: avanzamos en el tiempo para introducir ruido y retrocedemos en
    el tiempo para eliminarlo. Al leer otro trabajo en el que no
    queda claro el índice temporal, o al intentar implementar un algoritmo
    que presenta una ambigüedad similar, hay una forma de resolverla
    en todos los casos: el proceso de muestreo siempre debe tener un
    \textit{coeficiente positivo} en el término del filtro de ruido (o,
    equivalentemente, en la función de puntuación) al dar un paso.
\end{remark}

\begin{remark}\label{rem:flow_matching}
    Los procesos de difusión y eliminación de ruido, junto con los modelos asociados, ofrecen una
    forma de transformar la distribución de los datos hacia y desde una distribución de alta entropía
    (gaussiana o semejante a gaussiana). Por otro lado,
    puede resultar útil transformar la distribución de los datos hacia y desde una
    distribución objetivo arbitraria, que podría ser gaussiana, una
    mezcla gaussiana o incluso otro tipo de distribución de datos. Si
    tenemos este requisito, podemos utilizar el marco de \textit{flow matching}
    para lograrlo. Supongamos que la variable aleatoria objetivo es
    \(\vy\) y que nuestro flujo es
    \begin{equation}
        \vx_{t} = \alpha_{t}\vx + \sigma_{t}\vy, \qquad \forall t \in [0, 1].
    \end{equation}
    de tal manera que \(\alpha_{0} = \sigma_{1} = 1\) y \(\alpha_{1} =
    \sigma_{0} = 0\), de modo que \(\vx_{0} = \vx\) y \(\vx_{1} =
    \vy\). Nuestra estrategia no cambia en absoluto: podemos
    entrenar un filtro de ruido \(\bar{\vx}_{\theta}\) para eliminar el ruido de
    \(\vx_{t}\) y volver a obtener \(\vx\) (o, lo que es lo mismo, un predictor de «ruido»
    para predecir \(\vy\)) y luego aplicar nuestros algoritmos iterativos
    de eliminación de ruido para partir de una realización de \(\vy\) y
    obtener \(\vx\). Lo único que falta es la fórmula de Tweedie que
    relaciona la puntuación \(\nabla \log p_{t}\) del filtro de ruido con el
    filtro de ruido propiamente tal; dicha fórmula no existe en general.

    En la práctica, inspirándose en el concepto de aprender un así llamado
    \textit{mapa de transporte} entre pares de datos \(\vx\) y \(\vy\),
    los profesionales suelen querer entrenar un \textit{predictor de velocidad}
    \(\bar{\vv}_{\theta}\) para predecir \(\vv_{t} = \odv*{\vx_{t}}{t} =
    \alpha_{t}^{\prime}\vx + \sigma_{t}^{\prime}\vy\) en lugar de un
    filtro de ruido \(\bar{\vx}_{\theta}\), o más bien entrenar
    \(\bar{\vv}_{\theta}\) para predecir \(\Ex[\vv_{t} \mid \vx_{t}]\).
    En este caso, lo más sencillo es elegir los coeficientes
    \(\alpha_{t} = 1 - t\) y
    \(\sigma_{t} = t\) (para \(t \in [0, 1]\)), de modo que el objetivo de predicción
    sea \(\vv_{t} = \vy - \vx\) para cada instante \(t\). Dado el
    estimador \(\bar{\vv}_{\theta}\), el algoritmo de muestreo estándar
    recorre el campo hacia atrás en el tiempo; es decir, partiendo de una
    realización de \(\vy\), podemos actualizar iterativamente el iterado de la siguiente manera:
    \begin{equation}
        \hat{\vx}_{t_{\ell - 1}} = \hat{\vx}_{t_{\ell}} -
        \bar{\vv}_{\theta}(t_{\ell}, \hat{\vx}_{t_{\ell}}) \Delta t,
    \end{equation}
    donde \(\Delta t = t_{\ell} - t_{\ell - 1}\) es el intervalo de tiempo.
    Tenga en cuenta que la predicción de velocidad es algebraicamente equivalente a
    la eliminación de ruido (véase el \Cref{exercise:velocity_prediction}). Consulte las
    \Cref{sec:auto-encode-img-gen,sec:shape-and-pose} para ver aplicaciones prácticas
    del flow matching.
\end{remark}

\begin{remark}
    En la práctica, la difusión suele requerir muchos pasos y muchas evaluaciones
    de la función del filtro de ruido. Por lo tanto, es lógico preguntarse:
    «¿Podemos acelerar el proceso de muestreo utilizando un número menor
    de pasos y evaluaciones de la función?». La respuesta es afirmativa y a esos
    modelos se les suele llamar \textit{modelos de consistencia} o, de manera más
    descriptiva, \textit{modelos de pocos pasos}. La idea central consiste en utilizar un
    filtro de ruido que tome dos veces \(\bar{\vx}_{\theta}(s, t,
    \vx_{t})\) y estime \(\Ex[\vx_{s} \mid \vx_{t}]\), lo cual, en
    cierto sentido, es una tarea más sencilla, o bien utilizar un predictor de velocidad
    \(\bar{\vv}_{\theta}(s, t, \vx_{t})\) y estimar la velocidad media
    \(\Ex[(\vx_{t} - \vx_{s})/(t - s) \mid \vx_{t}]\). Estos
    filtros de eliminación de ruido tienen condiciones de consistencia
    adicionales, tales como \(\bar{\vx}_{\theta}(u, t,
    \bar{\vx}_{\theta}(s, u, \vxi)) = \bar{\vx}_{\theta}(s, t, \vxi)\),
    lo que justifica el uso de pérdidas de entrenamiento auxiliares. 
    Una vez que se aprende dicho predictor, el proceso de muestreo se vuelve 
    muy simple y eficiente, ya que solo requiere unos pocos pasos de la forma
    \begin{equation}
        \hat{\vx}_{s} = \bar{\vx}_{\theta}(s, t, \hat{\vx}_{t}) =
        \hat{\vx}_{t} - \bar{\vv}_{\theta}(s, t, \hat{\vx}_{t})\cdot(t - s).
    \end{equation}
    Entre las formulaciones relevantes se incluyen los mapas de flujo \cite{boffi2025build} y
    los flujos medios \cite{geng2025mean}.
\end{remark}

\section{Memorización, generalización y tasas de codificación}\label{sec:diffusion_memorization}
El cálculo que se presenta al final de la
\Cref{sub:intro_diffusion_denoising} parece indicar (en términos generales)
que, en la práctica, utilizar una red de tipo transformador es una
buena opción para entrenar o aproximar un filtro de ruido. Esto es
razonable, pero ¿qué problema hay en utilizar cualquier red neuronal
(como un perceptrón multicapa (MLP)) y simplemente ampliarla
hasta el infinito? Para responder a esta pregunta, consideremos una restricción
que hasta ahora hemos ignorado o pasado por alto en este capítulo: solo
disponemos de datos finitos con los que entrenar el filtro de ruido. De hecho, consideremos
el problema de entrenamiento \eqref{eq:parametric_denoising_loss},
excepto que esta vez suponemos que nuestros datos \(\vx\) se extraen de la
distribución empírica \(p^{\vX}\) sobre \(N\) muestras de entrenamiento
\(\vX = [\vx^{(1)}, \dots, \vx^{(N)}]\). Dado que \(p^{\vX}\) se define
en un subconjunto de dimensión \(0\) del espacio envolvente \(\R^{D}\),
tenemos que \(h(p^{\vX}) = -\infty\), por lo que, en cierto sentido, \(p^{\vX}\) es
la distribución más simple (es decir, la de menor entropía) que podría generar
nuestras observaciones \(\vX\). En resumen, ¡podría decirse que esta hipótesis \textit{tiene el
menor sesgo inductivo} respecto a la distribución de los datos! Muchos
comentarios populares sobre el aprendizaje profundo y los modelos generativos
quieren hacernos creer que, debido a esta propiedad, la mejor manera de resolver
el problema del modelado generativo es ampliar un filtro de ruido genérico
al máximo y aprovechar las ventajas correspondientes. Sin embargo, este consejo es (¡demostrablemente!) erróneo en este contexto.
Para entender por qué, observemos que la distribución empírica
\(p^{\vX}\) es, en realidad, una mezcla de distribuciones gaussianas con covarianzas nulas:
\begin{equation}
    p^{[\vx^{(1)}, \dots, \vx^{(N)}]} = \frac{1}{N}\sum_{i =
    1}^{N}\dNorm(\vx^{(i)}, \vzero).
\end{equation}
Por lo tanto, sabemos que la solución óptima (a nivel global) al
problema de eliminación de ruido \eqref{eq:parametric_denoising_loss} con \(\vx\)
extraído de la distribución empírica se expresa de la siguiente manera:
\begin{equation}\label{eq:parametric_denoising_loss_finite_sample}
    \min_{\theta \in \Theta}\frac{1}{N}\sum_{i =
    1}^{N}\Ex[w_{\tau}\norm{\vx^{(i)} - \bar{\vx}_{\theta}(\tau,
    \vx_{\tau}^{(i)})}_{2}^{2}]
\end{equation}
donde \(\vx_{t}^{(i)}\) es la versión con ruido de la muestra de entrenamiento
\(\vx^{(i)}\). El optimizador de esta pérdida sobre todos los
reductores de ruido (integrables en cuadrado) \(\bar{\vx}\) es el filtro de ruido gaussiano óptimo asociado
del \Cref{example:denoising_gaussian_mixture}, es decir,
\begin{equation}\label{eq:memorizing_denoiser}
    \bar{\vx}^{\star}(t, \vx_{t}) = \sum_{i =
    1}^{N}\frac{\exp({-\|\vx_{t} -
    \alpha_{t}\vx^{(i)}\|_{2}^{2}/(2\sigma_{t}^{2})})}{\sum_{j =
        1}^{N}\exp({-\|\vx_{t} -
    \alpha_{t}\vx^{(j)}\|_{2}^{2}/(2\sigma_{t}^{2})})}\vx^{(i)}.
\end{equation}
Esta es una combinación convexa de los datos \(\vx^{(i)}\) y los
coeficientes se vuelven más «agudos» (es decir, más cercanos a \(0\) o \(1\)) a medida que \(t
\to 0\). En particular, para valores pequeños de \(t\), el filtro de ruido es casi
exactamente igual a la muestra más cercana del conjunto de entrenamiento y la
actualización acerca la iteración al punto más cercano, de modo que al
final del muestreo, la iteración \(\hat{\vx}\) es (casi) igual a una
muestra de entrenamiento  \(\vx^{(i)}\). Insistamos de nuevo en este punto:
\textit{si utilizamos el filtro de ruido \eqref{eq:memorizing_denoiser} y un
buen muestreador, nuestras muestras regenerarán los datos de entrenamiento}. Esto
queda confirmado por nuestra garantía teórica sobre el muestreo
\Cref{thm:diffusion_sampler_convergence}, que, cuando se aplica en
este contexto, establece que si llevamos el número de pasos de muestreo a
\(\infty\) y utilizamos el programa de intervalos de tiempo prescrito, el modelo de difusión
asociado al filtro de ruido \eqref{eq:memorizing_denoiser}
siempre tomará muestras de la distribución empírica y, por lo tanto, siempre
reproducirá una muestra de entrenamiento. Este fenómeno constituye una aplicación especialmente estricta
de la \textit{memorización} en el modelado generativo.
\footnote{La memorización es un fenómeno mucho más amplio cuya
    caracterización exacta depende del dominio en el que se entrena el modelo
    (por ejemplo, texto, imágenes, vídeos, audio, etc.). Más adelante en esta sección
analizaremos medidas más heurísticas de la memorización.}

Dado que, en la práctica, solo disponemos de un conjunto de datos fijo y finito con el que podemos
entrenar nuestro filtro de ruido, el problema de aprendizaje que pretendemos resolver en este caso
es la pérdida de eliminación de ruido con muestra finita
\eqref{eq:parametric_denoising_loss_finite_sample} (en comparación con la
    pérdida de eliminación de ruido con muestra infinita \eqref{eq:parametric_denoising_loss},
por ejemplo). Por lo tanto, \textit{si tenemos un modelo arbitrariamente grande
    y lo optimizamos lo mejor posible, aprenderá a memorizar y
a regenerar el conjunto de entrenamiento}. En otras palabras, este modelo a gran escala
no es mucho más interesante que una simple base de datos con los datos de entrenamiento.

\begin{figure}
    \centering
    \includegraphics[width=\textwidth]{\toplevelprefix/chapters/chapter3/figs/euclidean_memorization_samples.png}
    \caption{\small\textbf{Muestras generadas (izquierda) a partir de un modelo de difusión
            entrenado con \(6{,}000\) muestras y (derecha) sus puntos más cercanos
        en el conjunto de datos de entrenamiento (en la norma euclidiana).} Observamos
        que las muestras generadas son, en esencia, equivalentes con una precisión de píxel
        a los datos de entrenamiento más similares. Ambos conjuntos de imágenes
    provienen de \cite{yoon2023diffusion}.}
    \label{fig:diffusion_memorization_euclidean_samples}
\end{figure}

\begin{figure}
    \centering
    \includegraphics[width=\textwidth]{\toplevelprefix/chapters/chapter3/figs/euclidean_memorization_ratio.png}
    \caption{\small\textbf{Memorización frente al tamaño del conjunto de datos para un
            modelo fijo, con (izquierda) 4.000 épocas de entrenamiento y (derecha)
        40.000 épocas de entrenamiento con conjuntos de datos más pequeños.} Observamos
        en el gráfico de la izquierda que los conjuntos de datos lo suficientemente grandes nunca se
        memorizarán (por ejemplo, de 10 000 a 50 000 muestras), mientras que los conjuntos de datos pequeños
        (de 1 000 a 5 000 muestras) se memorizarán tras
        un entrenamiento suficiente. Esto significa que, a medida que el conjunto de datos se vuelve
        más grande en relación con el modelo, la proporción de memorización al
        final del entrenamiento pasa de cerca de 1 a cerca de 0 con relativa
        rapidez, lo que indica la posible existencia de una transición de fase
        de la memorización a la generalización a medida que aumenta el tamaño
        del conjunto de datos. Figura tomada de
    \cite{gu2023memorization}.}
\end{figure}

El argumento anterior es válido para modelos de potencia arbitraria que puedan
ajustarse al (verdadero) filtro de ruido memorizador. Por lo tanto, es pertinente preguntarse
si esto se confirma realmente en la práctica: ¿los modelos muy grandes
(en relación con el conjunto de datos) aprenden a memorizar los datos de entrenamiento? Resulta
que esto es cierto desde el punto de vista empírico. Uno de los primeros estudios al respecto
fue \cite{yoon2023diffusion}. Para comprender sus resultados, definimos el
índice de memorización como la probabilidad de memorización:
\begin{equation}
    \mathrm{MemRatio}(\bar{\vx}) = \Pr_{\hat{\vx} \sim
    \mathrm{DM}(\bar{\vx})}[\hat{\vx}\ \text{is memorized}],
\end{equation}
donde \(\mathrm{DM}(\bar{\vx})\) es la distribución inducida por el
modelo de difusión que utiliza el filtro de ruido \(\bar{\vx}\), y la muestra
\(\hat{\vx}\) se memoriza utilizando el siguiente criterio:
\begin{equation}\label{eq:diffusion_euclidean_memorization_criterion}
    \hat{\vx}\ \text{is memorized}\ \iff \frac{\norm{\hat{\vx} -
    \vx_{(1)}(\hat{\vx})}_{2}}{\norm{\hat{\vx} -
    \vx_{(2)}(\hat{\vx})}_{2}} < \frac{1}{3},
\end{equation}
donde \(\vx_{(i)}(\hat{\vx})\) es el \(i\)-ésimo punto más cercano a
\(\hat{\vx}\) en el conjunto de datos de entrenamiento (en la norma euclidiana) y
\(1/3\) es una constante que, según se ha demostrado empíricamente, concuerda bien con
la percepción humana de la memorización. Nótese que esta definición se corresponde exactamente
con nuestra definición intuitiva anterior de «memorización», según la cual
cada uno de los puntos muestreados está mucho más cerca de una muestra de entrenamiento concreta que de todas
las demás muestras de entrenamiento. Demostraron que este tipo de memorización
se produce en la práctica cuando el modelo es mucho más grande (en términos del
número de parámetros) en comparación con los datos; la \Cref{fig:diffusion_memorization_euclidean_samples} muestra que,
en este caso, las muestras generadas son, en esencia, copias
perfectas, píxel a píxel, de las muestras de entrenamiento. Poco después, \cite{gu2023memorization}
llevó a cabo un estudio minucioso para caracterizar cómo diversos factores, tales como
las arquitecturas y el tiempo de entrenamiento, influyen en la memorización, demostrando, entre
otras cosas, que a medida que el tamaño del conjunto de datos aumenta en relación con un modelo fijo,
el índice de memorización cambia rápidamente de un valor cercano a 1 a uno cercano a 0,
es decir, el modelo deja de memorizar rápidamente. Esto
indica la posibilidad de una transición de fase entre la memorización
y la generalización a medida que el modelo se vuelve más pequeño en relación con los datos. 
\cite{buchanan2025edge} investiga la existencia y el mecanismo teórico
de esta transición de fase en datos sintéticos extraídos de un
modelo de mezcla gaussiana, demostrando que existe y proporcionando una caracterización parcial
del tamaño del modelo a partir del cual los modelos comienzan
o dejan de memorizar. Otros trabajos teóricos, como
\cite{niedoba2024towards,kamb2024analytic,george2025denoising,zhang2025generalization}
también intentan ofrecer una caracterización teórica de la transición de fase
en función del tamaño del modelo y del conjunto de datos en diferentes contextos.

\begin{figure}
    \centering
    \includegraphics[width=0.475\textwidth]{\toplevelprefix/chapters/chapter3/figs/generalizability.png}
    \hfill
    \includegraphics[width=0.45\textwidth]{\toplevelprefix/chapters/chapter3/figs/generalizability_loss.png}
    \caption{\small\textbf{La puntuación de generalización (izquierda) y
            la pérdida de entrenamiento (derecha) para un
        modelo fijo y un volumen de datos creciente}, extraído de
        \cite{pmlr-v235-zhang24cn}. La puntuación de generalización es \(1 - \)
        la fracción de muestras que se memorizan, y las diferentes
        líneas corresponden a modelos de eliminación de ruido de distintos tamaños
        (UNet-64 es más pequeño que UNet-128, que a su vez es más pequeño
        que UNet-256). Las cifras muestran que, a medida que el tamaño del conjunto de datos
        aumenta en relación con el tamaño del modelo, la pérdida aumenta
        significativamente y deja de producirse la memorización, lo que muestra que nuestra
    teoría e intuición se aplican a modelos reales.}
    \label{fig:diffusion_generalizability_loss}
\end{figure}

Con el fin de demostrar empíricamente la existencia de una transición de fase
entre la memorización y la generalización en datos reales,
\cite{pmlr-v235-zhang24cn} replanteó la definición de memorización
siguiendo el siguiente razonamiento. En el espacio de dimensiones muy elevadas de
las imágenes naturales, las distancias euclidianas entre puntos se vuelven mucho más
estrictas \cite{Wright-Ma-2022}, en el sentido de que, si dos imágenes no
coinciden a la perfección a nivel de píxeles, sus distancias euclidianas
se concentrarán de forma marcada en torno a un valor fijo (que varía en función de la
dimensión de las imágenes). Esto no tiene en cuenta los casos en los que ciertos
\textit{conceptos} (como los personajes de dibujos animados protegidos por derechos de autor) se
generan en posiciones novedosas, lo que dificulta enormemente establecer
criterios de memorización eficaces y prácticos basados en la distancia
euclidiana. Para solucionar esto, \cite{pmlr-v235-zhang24cn} utiliza una función concreta
\(f_{\mathrm{SSCD}}\) (una red neuronal preentrenada
denominada «detector de copias autosupervisado» \cite{pizzi2022self})
para mapear imágenes a un espacio euclidiano, y luego toma el producto interno
normalizado para obtener su similitud (coseno), de tal manera que las imágenes que
son similares debido a esta similitud coseno se consideran
conceptualmente similares en el sentido necesario para la memorización. Entonces,
\cite{pmlr-v235-zhang24cn} afirma que una muestra \(\hat{\vx}\) ha sido memorizada si su
similitud con cualquier punto de entrenamiento es \(> 3/5\), es decir,
\begin{equation}\label{eq:diffusion_sscd_memorization_criterion}
    \hat{\vx}\ \text{is memorized} \iff \max_{i \in [N]}
    \frac{f_{\mathrm{SSCD}}(\hat{\vx})^{\top}f_{\mathrm{SSCD}}(\vx^{(i)})}{\norm{f_{\mathrm{SSCD}}(\hat{\vx})}_{2}\norm{f_{\mathrm{SSCD}}(\vx^{(i)})}_{2}}
    > \frac{3}{5}
\end{equation}

La \Cref{fig:diffusion_generalizability_loss} muestra
el resultado de tomar numerosas muestras (alrededor de 10 000) de modelos de difusión
entrenados con conjuntos de datos de entrenamiento de distintos tamaños, en lo que respecta a
la estimación de la memorización y el cálculo de la pérdida de entrenamiento.
El resultado concuerda con lo expuesto anteriormente: cuando el modelo es
demasiado grande en relación con los datos de entrenamiento y la pérdida de entrenamiento es demasiado pequeña,
se produce memorización y las muestras reproducen los datos de entrenamiento. Así pues, nuestra
intuición se confirma en la práctica: \textit{¡en el caso de los modelos generativos de difusión, escalar
el modelo \underline{no} es lo único que se necesita!}

Entonces, la pregunta es: ¿qué tipo de filtro de ruido deberíamos usar? La clave
es elegir una arquitectura de red que pueda aproximar el verdadero
filtro de ruido (por ejemplo, el que corresponde a una distribución de rango bajo, como en
\eqref{eq:gmm_lowrank_denoiser}) sin aproximar el filtro de ruido memorizador
\eqref{eq:memorizing_denoiser}. Algunos trabajos han analizado esta
delgada línea y por qué los modelos de difusión modernos, que utilizan arquitecturas basadas en
Transformers y convoluciones, pueden memorizar y generalizar en
diferentes regímenes
\citep{kamb2024analytic,niedoba2024towards,zhang2025generalization}.

\begin{figure}
    \centering
    \includegraphics[width=0.4\textwidth]{\toplevelprefix/chapters/chapter3/figs/diffusion_perturbation.png}
    \caption{\small\textbf{Visualización del funcionamiento de un
        filtro de ruido bien entrenado} como una perturbación del filtro empírico en
        \textit{intervalos de tiempo cortos}, cuando el filtro con memoria actúa como una
    proyección sobre el punto más cercano del conjunto de entrenamiento.}
    \label{fig:diffusion_perturbation}
\end{figure}

Volvamos por fin a la premisa inicial de este capítulo, descrita
en la \Cref{sub:min_entropy}: encontrar la distribución de entropía mínima
que modela los datos es la única forma de aprender la distribución. Esperamos que
el análisis anterior haya demostrado que esta afirmación es falsa. Sin embargo,
una afirmación relacionada, también defendida en la \Cref{sub:min_entropy} debido a su
equivalencia en las situaciones que habíamos considerado entonces, podría seguir siendo
cierta: encontrar la distribución con la \textit{tasa de codificación} mínima es una
forma de aprender la distribución. Supongamos por un momento que \(\vx\)
tiene una densidad \(p\); la tasa de codificación (es decir, la entropía cruzada) de la
distribución empírica \(p^{\vX}\) con respecto a \(p\) es
\(\CE(p \mid p^{\vX}) = \infty\). Es evidente que esta tasa de codificación no es
mínima en ningún sentido; si utilizamos las mismas aproximaciones sucesivamente mejores
\(p^{n}\), \(p^{e}\) y \(\hat{q}\) analizadas en la
\Cref{sub:min_entropy}, se cumple que
\begin{equation}
    \infty = \CE(p \mmid p^{\vX}) > \CE(p \mmid p^{n}) > \CE(p \mmid
    p^{e}) > \CE(p \mmid \hat{q}) \approx h(p).
\end{equation}
Por lo tanto, la minimización de la tasa de codificación parece una forma sensata de conocer la
distribución y, al mismo tiempo, evitar este problema. Sin embargo, ahora surge un problema de
manejabilidad: ¿cómo podemos diseñar algoritmos que aprendan la
distribución minimizando la tasa de codificación y, al mismo tiempo, evitando los inconvenientes
asociados a una minimización excesiva de la entropía?

Una pista consiste en analizar qué ocurre con los modelos de difusión
entrenados con éxito que no memorizan. Al expresar el filtro de ruido entrenado
\(\bar{\vx}_{\theta}\) como una perturbación del filtro memorizador
\(\bar{\vx}^{\star}\), es decir,
\begin{equation}
    \bar{\vx}_{\theta}(t, \vx_{t}) = \bar{\vx}^{\star}(t, \vx_{t}) +
    \vdelta_{\theta}(t, \vx_{t})
\end{equation}
podemos hacernos una idea gráfica del funcionamiento de un filtro de eliminación de ruido bien entrenado,
como se muestra en la \Cref{fig:diffusion_perturbation}. A grandes rasgos,
el filtro de ruido entrenado calcula el filtro de memorización más una
perturbación. En un instante de tiempo \(t\) (quizás incluso en el último paso de la
eliminación de ruido), el filtro de ruido memorizador apunta completamente en la dirección
de una sola muestra de entrenamiento (la muestra de entrenamiento que se va a regenerar
en la iteración actual). La distribución que se está aprendiendo, o respecto a la cual se elimina el ruido,
es simplemente la distribución empírica con unas
esferas de radio \(\epsilon\) alrededor de cada punto,\footnote{Más formalmente, esta
    distribución perturbada tiene su soporte en (un subconjunto de) la unión de
    esferas de radio \(\epsilon\) centradas en cada muestra de entrenamiento. Esta imagen está
respaldada por trabajos anteriores, p. ej., \cite{song2025selective}.} donde
\(\epsilon\) es una cota superior de la magnitud de la perturbación,
visualizado en la \Cref{fig:diffusion_perturbation} por las esferas de color naranja 
pálido.\footnote{En dimensiones elevadas, se plantea la hipótesis de que un mejor
    modelo para estas perturbaciones no se describe mediante esferas, sino
    más bien mediante formas mucho más anisotrópicas y adaptadas a la geometría
\cite{farghly2025diffusion}, aunque la intuición sigue siendo similar.} En el caso de
modelos de difusión cada vez más grandes y bien entrenados, el
radio de perturbación \(\epsilon\) se reduce a \(0\) y el modelo
aprende la distribución empírica. Por otro lado, un valor de \(\epsilon\)
lo suficientemente grande como para que las bolas correspondientes a puntos diferentes se superpongan
o estén muy cerca puede implicar que el modelo de difusión sea generalizable.

Por lo tanto, nuestro modelo de difusión eficaz puede considerarse como una minimización de una
\textit{tasa de codificación con pérdida}, es decir, una tasa de codificación con una precisión
de \(\epsilon\) (de modo que se puedan ignorar las esferas de radio \(\epsilon\)
alrededor de cada punto). Esta noción de la tasa de codificación con pérdida resulta ser
lo que conviene minimizar, desde una perspectiva moral, y resultará
\textit{extremadamente} útil. Analizaremos sus aplicaciones en el
\Cref{ch:lossy-compression} y a lo largo del resto de este libro.

\section{Resumen y notas}
\paragraph{Mensajes clave.}
En este capítulo, hemos mostrado que, para obtener una distribución general
de baja dimensión, un enfoque computacional universal consiste
en partir de una distribución genérica y comprimirla hasta alcanzar
dicha distribución mediante la reducción (iterativa) del ruido o de la entropía. Como
vimos en el capítulo anterior, esto ocurre incluso en el caso de
las que podrían considerarse las distribuciones más simples de un único subespacio, como
el cálculo del PCA mediante el algoritmo de iteración de potencias. Como veremos en
los próximos capítulos, la función fundamental de las redes profundas modernas
no es ajustar distribuciones de datos o funciones arbitrarias (con ruido),
sino más bien implementar o aproximar esos operadores de compresión
o de eliminación de ruido (con pérdida). En capítulos posteriores derivaremos las formas precisas de esos
operadores y las arquitecturas profundas asociadas.

\paragraph{Más extensiones y referencias.}
El uso de la eliminación de ruido y la difusión para el muestreo tiene una larga trayectoria.
Probablemente, el primer trabajo que trata específicamente de un modelo de difusión sea
\cite{Sohl-Dickstein2015}, pero antes de este ya existían numerosos trabajos sobre
la eliminación del ruido como problema computacional y estadístico. El más
relevante de ellos es probablemente \cite{hyvarinen05a}, que utiliza explícitamente
la función de puntuación para eliminar el ruido y realizar un análisis de componentes 
independientes. Los trabajos posteriores más populares se basan, en esencia,
en la coocurrencia: \cite{ho2020denoising,song2019}. Desde entonces, miles
de artículos han partido de los modelos de difusión; retomamos este tema en el
\Cref{ch:consistent-representations}.

Muchos de estos trabajos utilizan un proceso estocástico distinto al de la
simple combinación lineal
\eqref{eq:gen_additive_gaussian_noise_model}. De hecho, todos los trabajos
mencionados anteriormente hacen hincapié en la necesidad de añadir ruido gaussiano \textit{independiente}
al inicio de cada paso del proceso hacia adelante. Los trabajos de carácter teórico 
utilizan el movimiento browniano o ecuaciones diferenciales estocásticas
para formular el proceso hacia adelante \cite{song2020score}. Sin embargo, dado que 
las combinaciones lineales de distribuciones gaussianas siguen dando lugar a 
distribuciones gaussianas, las \textit{distribuciones marginales} de
dichos procesos siguen adoptando la forma de
\eqref{eq:gen_additive_gaussian_noise_model}. La mayor parte de nuestro análisis
solo requiere que las distribuciones marginales sean las que son, por lo que
nuestro modelo, aunque sea demasiado simplista, resulta suficiente para casi todo.
El único caso en el que las distribuciones marginales no son suficientes es cuando derivamos
una expresión para \(\Ex[\vx_{s} \mid \vx_{t}]\) en términos de
\(\Ex[\vx \mid \vx_{t}]\). Los distintos procesos (de ruido) dan lugar a
diferentes expresiones de este tipo, que pueden utilizarse para el muestreo (y,
    por supuesto, existen otras formas de obtener muestreadores eficientes, como el
siempre popular muestreador DDPM). El proceso descrito en
\eqref{eq:gen_additive_gaussian_noise_model} es un auténtico
proceso estocástico cuya iteración de eliminación de ruido «natural» adopta la
forma del popular algoritmo DDIM
\cite{song2020denoising}. (Ni siquiera esta equivalencia es trivial;
citamos \cite{de2025distributional} como justificación.)

Además del trabajo teórico \citep{li2024d} tratado en la
\Cref{sub:sampling_denoising}, y la línea de trabajo en la que se basa,
que estudia la eficiencia de \textit{muestreo} de los modelos de difusión
cuando los datos tienen una estructura de baja dimensión, existe un amplio
conjunto de trabajos que estudia la eficiencia de \textit{entrenamiento} de
los modelos de difusión cuando los datos tienen una estructura de baja dimensión.
En concreto, \citet{chen2023score} y \citet{wang2024diffusion}
caracterizaron el error de aproximación y estimación de los filtros de ruido
cuando los datos pertenecen a una mezcla de gaussianas de bajo rango, demostrando
que el número de muestras de entrenamiento necesarias para aprender con precisión la
distribución varía proporcionalmente a la dimensión intrínseca de los
datos, en lugar de a la dimensión del entorno. Existe una considerable
\textit{labor metodológica} que intenta aprovechar la
estructura de baja dimensión de los datos para realizar diversas tareas con
modelos de difusión. Destacamos tres: la edición de imágenes
\citep{chen2024exploring}, la aplicación de marcas de agua \citep{li2024shallow} y
el «desaprendizaje» \citep{chen2025dual}, aunque esta lista no es exhaustiva.

\section{Ejercicios y extensiones}

\begin{exercise}
    Mostrar que \eqref{eq:optimal_denoiser} es la solución óptima del
    problema \eqref{eq:denoising_loss}.
\end{exercise}

\begin{exercise}\label{exercise:conditional_gaussian}
    Consideremos vectores aleatorios $\vx \in \bR^D$ y $\vy \in \bR^d$,
    de tal manera que el
    par $(\vx, \vy) \in \bR^{D + d}$ sea gaussiano conjuntamente. Esto significa que
    \begin{equation*}
        \begin{bmatrix}
            \vx \\
            \vy
        \end{bmatrix}
        \sim
        \cN \left(
            \begin{bmatrix}
                \vmu_{\vx} \\
                \vmu_{\vy}
            \end{bmatrix}
            ,
            \begin{bmatrix}
                \vSigma_{\vx} & \vSigma_{\vx\vy} \\
                \vSigma_{\vx\vy}^\top & \vSigma_{\vy}
            \end{bmatrix}
        \right),
    \end{equation*}
    donde los parámetros de media y covarianza vienen dados por
    \begin{equation*}
        \vmu_{\vx} = \bE[\vx],\quad \vmu_{\vy} = \bE[\vy],\quad
        \begin{bmatrix}
            \vSigma_{\vx} & \vSigma_{\vx\vy} \\
            \vSigma_{\vx\vy}^\top & \vSigma_{\vy}
        \end{bmatrix}
        =
        \bE\left[
            \begin{bmatrix}
                \vx - \bE[\vx] \\
                \vy - \bE[\vy]
            \end{bmatrix}
            \begin{bmatrix}
                \vx - \bE[\vx] \\
                \vy - \bE[\vy]
            \end{bmatrix}^\top
        \right]
    \end{equation*}
    Supongamos que $\vSigma_{\vy}$ es definida positiva (y, por lo tanto,
    invertible); entonces,
    la semidefinición positiva de la matriz de covarianza es equivalente
    a la condición del complemento de Schur
    $\vSigma_{\vx} - \vSigma_{\vx\vy}
    \vSigma_{\vy}^{-1}
    \vSigma_{\vx\vy}^\top \succeq \Zero$.

    En este ejercicio, demostraremos que la distribución condicional
    $p_{\vx \mid
    \vy}$ es gaussiana: es decir,
    \begin{equation}\label{eq:gaussian-conditional-eqn}
        p_{\vx \mid \vy} \sim \cN\left(
            \vmu_{\vx} + \vSigma_{\vx\vy} \vSigma_{\vy}^{-1} (\vy
            - \vmu_{\vy}),
            \vSigma_{\vx} - \vSigma_{\vx\vy} \vSigma_{\vy}^{-1}
            \vSigma_{\vx\vy}^{\top}
        \right).
    \end{equation}
    Una vía directa para demostrar este resultado consiste en manipular la relación definitoria de
    las densidades $p_{\vx, \vy} / p_{\vy}$. A continuación esbozamos un argumento algebraicamente conciso
    de este tipo.

    \begin{enumerate}
        \item Compruebe la siguiente identidad matricial para la covarianza:
            \begin{equation}\label{eq:gaussian-conditional-covariance-block}
                \begin{bmatrix}
                    \vSigma_{\vx} & \vSigma_{\vx\vy} \\
                    \vSigma_{\vx\vy}^\top & \vSigma_{\vy}
                \end{bmatrix}
                =
                \begin{bmatrix}
                    \vI_D & \vSigma_{\vx\vy}\vSigma_{\vy}^{-1} \\
                    \Zero & \vI_d
                \end{bmatrix}
                \begin{bmatrix}
                    \vSigma_{\vx} - \vSigma_{\vx\vy} \vSigma_{\vy}^{-1}
                    \vSigma_{\vx\vy}^{\top} & \Zero \\
                    \Zero & \vSigma_{\vy}
                \end{bmatrix}
                \begin{bmatrix}
                    \vI_D & \Zero\\
                    \vSigma_{\vy}^{-1}\vSigma_{\vx\vy}^\top & \vI_d
                \end{bmatrix}.
            \end{equation}
            Esta identidad se obtiene aplicando dos rondas de
            eliminación gaussiana
            en bloques a la matriz de covarianza.
        \item Partiendo de la identidad anterior, demuestre que
            \begin{equation}\label{eq:gaussian-conditional-covariance-inverse-block}
                \begin{bmatrix}
                    \vSigma_{\vx} & \vSigma_{\vx\vy} \\
                    \vSigma_{\vx\vy}^\top & \vSigma_{\vy}
                \end{bmatrix}^{-1}
                =
                \vP
                \begin{bmatrix}
                    \left(\vSigma_{\vx} - \vSigma_{\vx\vy}
                        \vSigma_{\vy}^{-1}
                    \vSigma_{\vx\vy}^{\top}\right)^{-1} & \Zero \\
                    \Zero & \vSigma_{\vy}^{-1}
                \end{bmatrix}
                \vP^\top,
            \end{equation}
            donde
            \begin{equation}
                \vP =
                \begin{bmatrix}
                    \vI_D & \Zero\\
                    -\vSigma_{\vy}^{-1}\vSigma_{\vx\vy}^\top & \vI_d
                \end{bmatrix},
            \end{equation}
            y siempre que las inversas pertinentes estén
            definidas.\footnote{En los casos en que el
                término del complemento de Schur no sea invertible, se mantiene el mismo
                resultado sustituyendo su
                inversa por la pseudoinversa de Moore-Penrose.
                En particular, la
                distribución condicional
                \eqref{eq:gaussian-conditional-eqn} se convierte en
            una distribución gaussiana degenerada.}
            Concluya que
            \begin{align}
                &
                \begin{bmatrix}
                    \vx-\vmu_{\vx} \\
                    \vy-\vmu_{\vy}
                \end{bmatrix}^{\top}
                \begin{bmatrix}
                    \vSigma_{\vx} & \vSigma_{\vx\vy} \\
                    \vSigma_{\vx\vy}^\top & \vSigma_{\vy}
                \end{bmatrix}^{-1}
                \begin{bmatrix}
                    \vx-\vmu_{\vx} \\
                    \vy-\vmu_{\vy}
                \end{bmatrix}
                \\
                &\qquad=
                \begin{bmatrix}
                    \vx - \vm \\
                    \vy - \vmu_{\vy}
                \end{bmatrix}^\top
                \begin{bmatrix}
                    \left(\vSigma_{\vx} - \vSigma_{\vx\vy}
                        \vSigma_{\vy}^{-1}
                    \vSigma_{\vx\vy}^{\top}\right)^{-1} & \Zero \\
                    \Zero & \vSigma_{\vy}^{-1}
                \end{bmatrix}
                \begin{bmatrix}
                    \vx - \vm \\
                    \vy - \vmu_{\vy}
                \end{bmatrix},
            \end{align}
            donde $\vm = \vmu_{\vx} +
            \vSigma_{\vx\vy}\vSigma_{\vy}^{-1}(\vy
            - \vmu_{\vy})$.
            %(\textit{Pista: Para simplificar las operaciones algebraicas, obsérvese
            % que la primera
            % y la última de las matrices del lado derecho de la
            % \Cref{eq:gaussian-conditional-covariance-block} son
            % transpuestas la una de la otra.})
        \item Dividiendo $p_{\vx, \vy} / p_{\vy}$, demuestre la 
            \Cref{eq:gaussian-conditional-eqn}. (\textit{Pista:
                    Utilizando las identidades anteriores,
                    solo debería ser necesario
                    un mínimo de álgebra.
                    Para la constante de normalización,
                    utilice la
                    \Cref{eq:gaussian-conditional-covariance-inverse-block}
                    para
            descomponer el determinante de manera similar.})

    \end{enumerate}

\end{exercise}

\begin{exercise}\label{exercise:sherman_morrison_woodbury_identity}
    Demuestre la identidad de Sherman-Morrison-Woodbury, es decir, para matrices
    \(\vA\), \(\vC\), \(\vU\), \(\vV\) tales que \(\vA\), \(\vC\) y
    \(\vA + \vU\vC\vV\) son invertibles,
    \begin{equation}
        (\vA + \vU\vC\vV)^{-1} = \vA^{-1} - \vA^{-1}\vU(\vC^{-1} +
        \vV\vA^{-1}\vU)^{-1}\vV\vA^{-1}.
    \end{equation}
\end{exercise}

\begin{exercise}\label{exercise:generalizing_results_to_different_noise_models}
    Derive de nuevo lo siguiente, suponiendo que \(\vx_{t}\) sigue el modelo de ruido generalizado
    \eqref{eq:gen_additive_gaussian_noise_model}.
    \begin{itemize}
        \item La fórmula de Tweedie: \eqref{eq:gen_tweedie}.
        \item La iteración del DDIM: \eqref{eq:gen_denoising_iteration}.
        \item El filtro de ruido óptimo de Bayes para un modelo de mezcla gaussiana:
            \eqref{eq:gen_gmm_bayes_optimal_denoiser}.
    \end{itemize}
\end{exercise}

\begin{exercise}\label{exercise:implement_denoising_processes}
    \phantom{}
    \begin{enumerate}
        \item Implemente las fórmulas derivadas en el
            \Cref{exercise:generalizing_results_to_different_noise_models},
            creando un muestreador para mezclas gaussianas.
        \item Reproduzca la \Cref{fig:ve_forward_denoising} y la
            \Cref{fig:vp_gmm_denoising}.
        \item A continuación, presentamos \textit{Flow
            Matching (FM)} (véase la \Cref{rem:flow_matching}), de la siguiente manera:
            \begin{equation}
                \alpha_{t} = 1 - t, \qquad \sigma_{t} = t.
            \end{equation}
            Implemente este proceso utilizando el mismo marco y pruébelo
            para el muestreo en espacios de alta dimensión. ¿Qué proceso ofrece resultados mejores
            o más estables?
    \end{enumerate}
\end{exercise}

\begin{exercise}\label{exercise:velocity_prediction}
    Consideremos la interpolación, que nos resulta algo familiar
    \begin{equation}
        \vx_{t} = \alpha_{t}\vx + \sigma_{t}\vy, \qquad \forall t
        \in [0, T].
    \end{equation}
    donde \(\vy\) es una variable aleatoria (no necesariamente gaussiana).
    Definimos la \textit{velocidad} como
    \begin{equation}
        \vv_{t} = \odv*{\vx_{t}}{t} = \alpha_{t}^{\prime}\vx +
        \sigma_{t}^{\prime}\vy, \qquad \forall t \in [0, T].
    \end{equation}
    Calcule una relación algebraica que exprese \(\Ex[\vv_{t} \mid
    \vx_{t}]\) en términos de \(\Ex[\vx \mid \vx_{t}]\) y viceversa,
    demostrando que la predicción de la velocidad es algebraicamente
    equivalente a la eliminación de ruido y, por lo tanto, también a la predicción del ruido
    (véase la \Cref{rem:flow_matching}).
\end{exercise}

\end{document}
